\documentclass[preprint,12pt]{elsarticle}
\usepackage{amsmath,amssymb,amsthm,booktabs,graphicx,siunitx,xcolor,array,float}
\PassOptionsToPackage{hyphens}{url}
\usepackage[T1]{fontenc}
\usepackage{lmodern}
\usepackage{microtype}
\usepackage[pdfusetitle=false]{hyperref}
\setlength{\emergencystretch}{4em}
\hbadness=5000
\hyphenpenalty=700
\tolerance=1500
\hyphenation{transfers consistent consistency Supplementary unannounced Newton Gronwall Lipschitz Runge Kutta Riccati MATPOWER Jacobian Willems Kunkel Mehrmann Messenger Hamiltonian Lyapunov Bergen}
\graphicspath{{figures/}}
\journal{Communications in Nonlinear Science and Numerical Simulation}
\makeatletter\def\ps@pprintTitle{\let\@oddhead\@empty\let\@evenhead\@empty\def\@oddfoot{\scriptsize\itshape Preprint submitted to Communications in Nonlinear Science and Numerical Simulation\hfill October 6, 2026}\let\@evenfoot\@oddfoot}\makeatother
\newtheorem{proposition}{Proposition}
\newtheorem{remark}{Remark}
\newcommand{\dd}{\mathrm{d}}
\newcommand{\RR}{\mathbb{R}}
\newcommand{\repodoi}{[repository DOI to be inserted before submission]}
\newcommand{\chainA}{1.73}
\newcommand{\chainB}{2.03}
\newcommand{\chainC}{5.13}
\newcommand{\chainD}{9.84}
\hypersetup{pdftitle={Structure-preserving residual learning and predictive control of nonlinear descriptor power-network models},pdfauthor={Sri Venkata Durga Sudarsan Madhyannapu},colorlinks=true,linkcolor=blue!50!black,citecolor=blue!50!black,urlcolor=blue!50!black}

\begin{document}
\begin{frontmatter}
\title{Structure-Preserving Residual Learning and Predictive Control of Nonlinear Descriptor Power-Network Models}
\author[a]{Sri Venkata Durga Sudarsan Madhyannapu\corref{cor1}}
\ead{msvdsudarsan@gmail.com}
\cortext[cor1]{Corresponding author. ORCID: 0009-0001-2126-6428}
\affiliation[a]{organization={Department of Mathematics, School of Sciences, Humanities and Management, Dr.~RVR~NRI~Institute of Technology (Deemed to be University)},
addressline={Vijayawada Rural},
city={Vijayawada}, postcode={521212}, state={Andhra Pradesh}, country={India}}

\begin{abstract}
A learned correction to the differential part of a differential-algebraic (descriptor) model is useful only if the corrected predictor stays on the algebraic constraint manifold. This paper asks when structure-preserving residual learning improves the prediction and control of a nonlinear descriptor model, when it stops helping, and what to do then. The known physics and algebraic law are kept; only the missing differential term is identified, in weak form, from noisy trajectories. A Newton projection returns the algebraic variables to the manifold, and a split-conformal margin shifts the penalty thresholds of a receding-horizon controller. Matched variants separate these ingredients on a scalar benchmark and on reduced network realizations of MATPOWER case9 and case39. An eigenvalue bound for the algebraic Jacobian gives a Lipschitz constant, a Gronwall bound on residual-error growth and an exact gradient that brings the case39 step time to about 24 ms. The corrected predictor reduces the one-second state error to 0.17 (case9) and 0.24 (case39) of the physics-only value. In the reduced models a validity atlas shows that the benefit varies much more strongly with the dispatch error than with the loading. A trust gate detects an unannounced dispatch step within 0.3 s and falls back to the physics model; a reconciliation step then re-estimates a common dispatch scale to within 0.5\% and restores the correction. The margin, the gate and the reconciliation are heuristics without guarantees, and the models are not electromagnetic-transient simulations.
\end{abstract}
\begin{keyword}
descriptor systems \sep differential-algebraic equations \sep model predictive control \sep residual learning \sep conformal prediction \sep power networks \sep system identification
\end{keyword}
\end{frontmatter}

\section{Introduction}
\label{sec:intro}
Network models in power systems, constrained mechanics and process engineering couple differential evolution with algebraic balance laws. In semi-explicit form they read
\begin{equation}
\dot x = f(x,z,u),\qquad 0 = g(x,z,u),
\label{eq:dae}
\end{equation}
with differential variables $x$, algebraic variables $z$ and inputs $u$. The algebraic relation is more than bookkeeping. It cuts out the admissible set $\mathcal M=\{(x,z,u):g(x,z,u)=0\}$, and a trajectory that leaves $\mathcal M$ is not a trajectory of the model. When $\partial g/\partial z$ is nonsingular the system has index one and $z$ can in principle be eliminated, but in a large network the elimination is rarely carried out explicitly and the constraint is imposed numerically.

Reduced physical models leave out effects such as saliency, nonlinear friction or load dynamics, and a data-driven correction is the obvious remedy. The correction, however, enters the differential equations only, so nothing in it protects the constraint. If the algebraic variables are advanced inside a predictor with a cheap approximation, the learned term can improve one-step accuracy while the constraint residual grows along the horizon. A black-box model of $(x,u)\mapsto\dot x$ is no better placed: it has to rediscover what the physics already supplies, and it does not respect $g=0$ by construction. Here the approach is the reverse. The known structure ($f$ and $g$) is kept, learning is asked to supply only what is missing, and the algebraic variables are restored by a Newton projection at every prediction stage. Because $z=h(x,u)$ on $\mathcal M$, the learned residual is in effect a function on the manifold, and Section~\ref{sec:theory} shows how the regularity of $\partial g/\partial z$ governs the growth of its error.

The learned correction is not assumed to be valid everywhere. The framework detects when its operating assumptions fail, gives priority to the known physics for a while, and, when the mismatch can be reconciled, restores the learned correction. The question of the paper follows from this. \emph{When does structure-preserving residual learning improve the prediction and control of a nonlinear descriptor model, and when does it fail?} To answer it, the residual, the projection onto $\mathcal M$ and an empirical penalty margin are added one at a time (Fig.~\ref{fig:pipe}, Table~\ref{tab:variants}). Each step of the ladder C0, A0, A1, A2, A3 then isolates one ingredient, and its effect can be read from prediction accuracy, constraint consistency, closed-loop behaviour and computing cost.

The contribution is not a new learner and not a new MPC solver. It is a structure-preserving combination of residual identification, projection onto the algebraic manifold, finite-sample calibration of a prediction margin and receding-horizon control for a nonlinear descriptor model, together with an explicit account of the operating conditions and the computing cost at which the correction stops paying off.

\begin{figure*}[!t]\centering
\includegraphics[width=\linewidth]{fig_pipeline.pdf}
\caption{Pipeline and matched variants. The residual is learned for the differential part only, the algebraic variables are restored by Newton projection, and an empirical margin on the multi-step error shifts the penalty thresholds of the controller. C0--A3 are defined in Table~\ref{tab:variants}. The drawing shows no data.}
\label{fig:pipe}\end{figure*}

The paper makes five contributions. They follow one line of thought: known physics, learning of the missing differential term, preservation of the algebraic manifold, a statement of where learning helps, a gate that withdraws the correction when its operating assumptions fail, and a step that tries to restore them.
\begin{enumerate}
\item Two reproducible nonlinear descriptor benchmarks whose physics-only models are deliberately mismatched with the plant: a scalar benchmark with one algebraic voltage-like variable, and a network realization built from the published MATPOWER case9 and case39 data with all bus angles as algebraic variables.
\item A pipeline that identifies the residual from noisy data in weak form, with disjoint training, validation, calibration and test trajectories, and turns held-out multi-step errors, collected per trajectory under a static law and under the receding-horizon controller, into an empirical envelope (a split-conformal quantile over trajectories) that shifts the penalty thresholds of the controller. Matched variants isolate the residual, the projection and the tightening. The envelope is a calibrated heuristic, its coverage is reported, and it is not a safety guarantee.
\item Three checkable statements, read as one chain: the scalar algebraic equation is index one globally; the algebraic Jacobian of the network realization has an explicit eigenvalue lower bound whenever angle differences stay within a stated limit; and the same bound gives a Lipschitz constant for the reduced dynamics, which sets the growth rate in a Gronwall bound for the propagation of the residual error.
\item A map of where the correction stops helping. The residual transfers across loadings when the dispatch is known but not under an unannounced shift of the dispatch; the penalty tightening is inactive in the network experiments; and the network residual is not sparse. The case39 controller with finite-difference gradients is slower than the sampling period; an exact gradient of the projected predictor (Section~\ref{sec:grad}) removes this limit without changing the closed-loop result. Fidelity, protection and computing cost are reported side by side (Section~\ref{sec:regime}).
\item A validity atlas over load scale and dispatch error (Section~\ref{sec:atlas}), a transparent residual trust gate that falls back to the physics model when the corrected predictor stops explaining the measurements, and a reconciliation step that re-estimates the dispatch and restores the correction (Sections~\ref{sec:gate} and~\ref{sec:gate_results}). The mechanism is a heuristic with no guarantee; it keeps the Newton projection active in every state.
\end{enumerate}
The three test beds answer different questions. The scalar benchmark is small enough to be analysed in closed form (a global index-one property, a basin estimate, an exactly known residual), so it shows the mechanism transparently. Case9 asks whether the mechanism survives the algebraic structure of a network. Case39 adds scale, and with it the question of what the more faithful predictor costs in computing time.

The work does not claim electromagnetic-transient validation, a closed-loop stability certificate or a worst-case robustness guarantee. The margin is an empirical quantity, and the title makes no claim of robustness.

\section{Related work}
\label{sec:related}
Differential-algebraic equations are treated in depth in \cite{brenan,kunkel}; regularity and index questions for nonlinear systems are discussed in \cite{scholz} and stabilization of nonlinear DAEs in \cite{liuho}. Linear descriptor systems have recently been combined with data-driven predictive control, including a nine-bus power-network case study \cite{schmitz1,schmitz2}. Those results use a discrete-time linear descriptor representation and a behavioural (Willems-type) predictor; the present work keeps a nonlinear physical model and learns a bounded correction to it.

Sparse regression of governing equations \cite{brunton} has been used inside predictive control in the low-data regime \cite{kaiser}, and the weak (integral) formulation \cite{wsindy} avoids differentiating noisy data. Learning-based predictive control is reviewed in \cite{hewing}. Constraint tightening against bounded disturbances is classical \cite{mayne}, and probabilistic tightening has been developed for stochastic predictors \cite{lorenzen}. Conformal prediction \cite{vovk,angelopoulos} supplies distribution-free finite-sample quantiles under exchangeability; here it is used to calibrate multi-step prediction-error envelopes, with the caveat that closed-loop windows are not exchangeable (Section~\ref{sec:limits}). Weighted and adaptive variants of conformal prediction address covariate shift and drift in the data \cite{tibshirani,barber,gibbs}; the gate of Section~\ref{sec:gate} does something simpler, because it does not recalibrate an interval but attenuates the learned correction when the predictor stops explaining the measurements. Physics-informed learning of DAEs on power networks was proposed in \cite{moya}, and recent work on structure-preserving reduction of nonlinear power-network DAEs is reported in \cite{nadeem,banagaaya}. The network-preserving formulation used below, with generator internal buses and algebraic load-bus balance, goes back to \cite{bergenhill}; see also \cite{dorfler,kundur}. Reviews of data-driven power-system dynamics emphasize data quality, operating-condition variation and model transparency \cite{huang}, which is the framing adopted for the experiments on parameter and operating-point shift in Section~\ref{sec:sens}. Recent contributions to this journal treat delay-dependent stability of multi-area load frequency control with electric vehicles \cite{xiong} and structure-preserving neural networks for the Fermi--Pasta--Ulam--Tsingou lattice \cite{fan}; the first concerns delay-dependent stability of multi-area control and the second a structure-preserving network for a nonlinear lattice, whereas the present paper studies learned corrections of nonlinear differential-algebraic power-network models inside a predictive controller. The author's earlier work analysed the breakdown of observability in singular bilinear periodic matrix differential systems with a Melnikov-based argument \cite{madh}, and developed an adjoint-free Melnikov--Lyapunov diagnostic for transition detection in slow--fast chemical reaction networks \cite{madh2}. Both concern structural or diagnostic properties of other system classes; they contain no learning, prediction or control, and none of their results is reused here. The present paper keeps the interest in singular structure and moves from analysis and diagnosis to prediction and control, asking how an algebraic constraint shapes what can be learned and predicted for a network model.

\section{Problem setting and benchmarks}
\label{sec:setting}
\subsection{Residual-augmented descriptor model}
The plant is a semi-explicit DAE
\begin{equation}
\dot x = f(x,z,u) + d(x,z,u) + w(t),\qquad 0=g(x,z,u),
\label{eq:plant}
\end{equation}
in which the physics model $f$ and the algebraic law $g$ are known, $d$ is an unknown residual and $w(t)$ an exogenous disturbance, measured in units of the state derivative. In the scalar benchmark $w=(0,p(t)/M)$, see \eqref{eq:s2}. The network realization has no additive term of this kind: its load fluctuation enters the algebraic law, which then depends on $P_d(t)$ (Section~\ref{sec:network}). The learned predictor is $\dot x=f+\hat d$, $0=g$. When $\partial g/\partial z$ is nonsingular, $z=h(x,u)$ on the manifold $\mathcal M=\{g=0\}$ and $d$ restricted to $\mathcal M$ is a function of $(x,u)$; the residual is therefore learned on the manifold, and the predictor leaves it only if $z$ is not recomputed from $g=0$. The method is organised in three layers: the singular structure $(f,g)$, which is kept; the identification of $d$ on $\mathcal M$; and the use of the identified residual, with its empirical margin, inside a predictor and a controller. Throughout, $g$ is treated as known because it expresses a conservation law; mismatch is placed in the differential part. This is an assumption of the study, not a property of real networks, where line parameters are also uncertain (Section~\ref{sec:sens}).

\subsection{Scalar benchmark}
\label{sec:scalar}
The scalar benchmark has the angle-like state $\delta$, the frequency-like state $\omega$ and one algebraic voltage-like variable $v$:
\begin{align}
\dot\delta &= \omega, \label{eq:s1}\\
M\dot\omega &= P - k_v\,v\sin\delta - D\omega + u - c_3\delta^3 + p(t), \label{eq:s2}\\
0 &= g(\delta,v,u) = v + a\sin\delta + c_b v^3 - 1 - k_u u, \label{eq:s3}
\end{align}
with $M=0.55$, $P=1.10$, $k_v=0.85$, $D=0.18$, $c_3=0.10$, $a=0.12$, $c_b=0.08$, $k_u=0.04$, input bound $|u|\le 0.8$ and $p(t)=A\sin(\Omega t+\phi)$; in the notation of \eqref{eq:plant} the disturbance on the frequency state is $p(t)/M$. The physics-only model uses $k_v=0.80$ and $D=0.14$ and omits the cubic term, so that the residual $d$ contains a parametric part and a structural part. It is a power-network-inspired toy model, not a calibrated machine model.

\subsection{Network realization of MATPOWER case9 and case39}
\label{sec:network}
Topology, line reactances $x_l$, bus loads, generator buses, active-power dispatch and generator ratings are read from the official MATPOWER files \cite{matpower,matpowerdata}. Each generator $i$ at bus $\sigma(i)$ has the internal angle $\delta_i$ and speed deviation $\omega_i$; all $n$ bus angles $\theta_j$ are algebraic variables of a lossless network with unit voltage magnitudes:
\begin{align}
\dot\delta_i &= \omega_i,\nonumber\\ M_i\dot\omega_i &= P_{m,i} + u_i - D_i\omega_i - b_i\sin(\delta_i-\theta_{\sigma(i)}), \label{eq:n1}\\
0 &= g_j(\delta,\theta) = \sum_{i:\sigma(i)=j} b_i\sin(\delta_i-\theta_j) - P_{d,j} \nonumber\\ &\quad - \sum_{l\sim j} b_l\sin(\theta_j-\theta_{k(l,j)}), \label{eq:n2}
\end{align}
where $b_l=1/x_l$ is the line susceptance and $b_i=1/x'_{d}$ the generator-to-bus susceptance. The dispatch is scaled uniformly so that $\sum_i P_{m,i}=\sum_j P_{d,j}$. MATPOWER does not provide inertia, damping or transient reactance, so these are declared assumptions: $H=4$\,s on the generator rating, $x'_d=0.3$ per unit on that rating, $D_i/M_i=4\ \mathrm{s}^{-1}$ and an input bound $|u_i|\le 0.1\,S_i$ (\ref{app:params}). The plant contains three effects that the physics model does not: a saliency term $\kappa b_i\sin(2(\delta_i-\theta_{\sigma(i)}))$ with $\kappa=0.08$ in the electrical power, a saturating friction $0.15\,D_i\tanh(2\omega_i)$ and a sinusoidal load fluctuation $P_d(t)=P_d(1+\varepsilon\sin(\Omega t+\phi))$. The physics model has damping $0.8D_i$ and none of these effects, and its predictions hold the load at the value it has at the start of the prediction window. The case9 realization has $n=9$ buses and $3$ generators ($6$ differential states, $9$ algebraic variables); case39 has $39$ buses and $10$ generators ($20$ differential states, $39$ algebraic variables). Angles are reported relative to the centre of inertia.

The realization is a reduced swing model that keeps the network. MATPOWER files contain steady-state power-flow data, so exciter, governor, load-dynamic and voltage-magnitude effects are absent by construction, and no statement about electromagnetic-transient behaviour is made.

\section{Method}
\label{sec:method}
\subsection{Residual identification in weak form}
\label{sec:ident}
For each generator, the residual acceleration is expressed in a library of $13$ features of the rotor angle $\psi_i=\delta_i-\theta_{\sigma(i)}$, the speed and the input: $\sin k\psi_i$, $\cos k\psi_i$ ($k=1,2,3$), $\omega_i$, $\omega_i^2$, $\omega_i^3$, $u_i$, $u_i\omega_i$, a constant and $\psi_i$. The scalar benchmark uses all monomials in $(\delta,\omega,u)$ up to degree $4$. Measured trajectories are noisy ($\sigma=10^{-3}$ on every state), sampled at step $h$ with a zero-order-hold input, and the regression target is the trapezoidal weak form
\begin{equation}
\frac{\omega_{k+1}-\omega_k}{h} - \tfrac12\bigl(a^{\mathrm{ph}}_k + a^{\mathrm{ph}}_{k+1}\bigr) = \tfrac12\bigl(\Phi_k+\Phi_{k+1}\bigr)c ,
\label{eq:weak}
\end{equation}
where $a^{\mathrm{ph}}$ is the physics-model acceleration with the algebraic variables evaluated from $g=0$ at the measured states, and $\Phi$ is the feature vector. This form needs no numerical differentiation of the frequency signal and is insensitive to the input jumps of a sampled controller. Coefficients are obtained by sequentially thresholded ridge regression on RMS-normalized features. The ridge parameter and the threshold are chosen on the validation set by a parsimony rule: among all candidates whose validation RMSE lies within $5\%$ of the best, the one with the fewest active terms is kept. The result is called a thresholded regression model; as Section~\ref{sec:results} shows, it is parsimonious on the scalar benchmark and only moderately so on the networks. The fitted correction is clipped at $1.5$ times the largest training residual, which enforces the bounded-residual assumption outside the training region.

Data are split by trajectory into training, validation, calibration and test sets (scalar benchmark: $24/8/10/16$ trajectories; network: $30/8/8/16$ per loading). Excitation comes from random initial perturbations, random input dither and, in the network case, load fluctuations of amplitude up to $4\%$. The test set is generated without measurement noise so that the exact residual is available for evaluation; it is never used for fitting, selection or calibration.

\paragraph{Measurement model.} In the network experiments the generator angles $\delta_i$ and speeds $\omega_i$ are measured with additive Gaussian noise ($\sigma=10^{-3}$); the bus angles are not measured separately but are reconstructed from the measured $\delta$ by solving $g=0$. The algebraic variables are therefore exact functions of the noisy differential states, and the study does not test the case in which the algebraic variables are measured with independent error.

\subsection{Algebraic projection}
Two ways of propagating the algebraic variables inside the predictor are compared. In the tangent update, $z$ is advanced with the first-order sensitivity from the last consistent point, $z_{+}=z-(\partial_z g)^{-1}\partial_x g\,\Delta x$, which is the cheap option when no corrector is available. In the projection, a Newton corrector solves $g=0$ at every Runge--Kutta stage and at the end of the step, with tolerance $10^{-10}$. The same stage structure is used in both, so any difference comes from the treatment of $g$ alone. The plant itself is integrated with $g=0$ solved to $10^{-11}$.

\subsection{Prediction-error envelope and penalty tightening}
\label{sec:tube}
The envelope is built from trajectories rather than from individual windows, because windows cut from one trajectory overlap and are strongly dependent. For calibration trajectory $m=1,\dots,n$ let $s_k^{(m)}=\max_{\text{windows}}|e_k^{(m)}|$, where $e_k^{(m)}$ is the error at prediction step $k$ of a constrained output (the scalar benchmark: $\delta$ and $\omega$; the network: the largest line-angle difference), taken over all windows of that trajectory. The envelope is the split-conformal quantile across trajectories,
\begin{equation}
\begin{split}\varepsilon_k ={}&\lceil (n+1)(1-\alpha)\rceil\text{-th smallest}\\ &\text{of } \{s_k^{(m)}\}_{m=1}^{n},\ \alpha=0.05,\end{split}
\label{eq:conf}
\end{equation}
and for $\alpha=0.05$ it is simply the largest score when $n\le38$. The calibration set mixes trajectories under the static law and under the receding-horizon controller without tightening (A2), generated from the disturbance distribution used at deployment, so that the calibration policy matches the deployment policy. The controller then replaces each state limit $\bar y$ at prediction step $k$ by $\bar y-\varepsilon_k$. Because the limits enter the cost as quadratic penalties (Section~\ref{sec:mpc}), this moves the threshold at which the penalty starts; it does not enforce a hard state constraint, and it is called penalty tightening here. The envelope is empirical. Its nominal level is not a guarantee: the number of calibration trajectories is small, trajectories under different policies are only approximately exchangeable, and the coverage actually obtained on independent trajectories is reported in Section~\ref{sec:results}. For the one-step residual error of the scalar benchmark the same construction is applied to maxima over non-overlapping blocks of $20$ samples ($1$\,s) and gives the bound $\rho_{\rm blk}$; a pointwise quantile $\rho_{\rm pt}$ is reported next to it for comparison.

Calibration is included for a reason. A small regression error on one-step residual targets says little about the error of a multi-step prediction, because the predictor is driven by its own output and the error compounds along the horizon. The envelope is therefore measured on the quantity the controller actually uses, the multi-step error of the projected predictor, and is not inferred from the fit. Four objects are kept apart below: the analytical bound of Proposition~\ref{prop:gronwall}, which is a worst-case statement under assumptions; the empirical envelope; the conformal quantile that calibrates it; and the closed-loop performance that results. Only the last is a measured outcome of control, and none of the first three is a guarantee for the closed loop.

\subsection{Receding-horizon controller and baselines}
\label{sec:mpc}
The controller solves, every $T_s=0.1$\,s, a finite-horizon problem with sampled input, the corrected predictor and a quadratic stage cost on the tracking error, the frequency and the input and its increments; constraints are imposed as quadratic penalties and the input bound as a box. The scalar benchmark uses $N=12$ steps, a terminal cost from the discrete algebraic Riccati equation of the model linearized at the set point, and the limits $\delta\le 0.66$, $|\omega|\le 0.60$. The network uses $N=5$ steps without a terminal cost, so no feasibility or stability property is claimed, and the limits $|\theta_j-\theta_k|\le 30^\circ$ on every line and $|\delta_i-\theta_{\sigma(i)}|\le 30^\circ+0.35$\,rad. The problem is solved by \mbox{L-BFGS-B} with a warm start from the shifted previous solution (at most $60$ iterations for the scalar case, $12$ for the network). The gradients are finite differences, except for the projected network predictor (A2, A3), where the exact gradient of Section~\ref{sec:grad} is used. Table~\ref{tab:variants} lists the variants. C0 is a static law: for the scalar benchmark $u=u_{ff}-1.25(\delta-\delta_{ref})-0.52\omega-0.06\,\mathrm{sgn}(\omega)\min\{|\omega|,1\}$ with $u_{ff}$ from the physics-model equilibrium, saturated at $\pm0.8$; for the network damping injection $u_i=-D_i\omega_i$, saturated.

\begin{table*}[!t]
\centering\small
\caption{Variants compared. A0--A2 differ only in the predictor; A3 adds the penalty tightening of Section~\ref{sec:tube}; G adds the gate and the reconciliation to A2. Each step isolates one structural ingredient.}
\label{tab:variants}
\begin{tabular}{l>{\raggedright\arraybackslash}p{3.3cm}ccc>{\raggedright\arraybackslash}p{3.3cm}}
\toprule
Label & Description & Residual & Proj. & Tight. & Ingredient isolated\\
\midrule
C0 & static feedback & -- & -- & -- & reference without a predictor\\
A0 & physics-only predictor, tangent update & no & no & no & mismatched physics\\
A1 & A0 with learned residual & yes & no & no & missing differential physics\\
A2 & A1 with Newton projection & yes & yes & no & preservation of the manifold\\
A3 & A2 with penalty tightening & yes & yes & yes & empirical error margin\\
G & A2 with trust gate and dispatch reconciliation (Section~\ref{sec:gate}) & yes & yes & no & trust in the correction\\
\bottomrule
\end{tabular}
\end{table*}

\subsection{Exact gradient for the projected predictor}
\label{sec:grad}
With the Newton projection the algebraic variables of the predictor are, to solver tolerance, the function $\theta(\delta)$ defined by $g=0$, and Proposition~\ref{prop:lip} gives $\partial\theta/\partial\delta=J^{-1}EC$. The projected predictor is therefore an ordinary Runge--Kutta integration of the reduced model $x=(\delta,\omega)$, and the cost of Section~\ref{sec:mpc} can be differentiated exactly with respect to the input sequence. In every stage the Jacobians of the vector field $F$ are
\begin{equation}
\begin{split}F_x&=\begin{pmatrix}0&I\\ S_\psi\,(I-E^\top J^{-1}EC)&S_\omega\end{pmatrix},\\ F_u&=\begin{pmatrix}0\\ S_u\end{pmatrix},\end{split}
\label{eq:stagejac}
\end{equation}
where the diagonal matrices $S_\psi$, $S_\omega$ and $S_u$ hold the derivatives of the generator acceleration (physics model plus learned residual, including its clipping) with respect to $\psi_i$, $\omega_i$ and $u_i$. Chaining the four stages gives the sensitivities $\partial x^{+}/\partial x$ and $\partial x^{+}/\partial u$ of one step, and composing the sub-steps gives one pair of matrices per control interval. The penalty terms also depend on $\delta$ through $\theta(\delta)$, which adds the term $(J^{-1}EC)^\top$ times the penalty gradient with respect to $\theta$. A backward recursion over the horizon then gives the gradient with respect to all inputs. Each stage needs one factorization of $J$ and a solve with $n_g$ right-hand sides, so the cost is about that of two predictor rollouts, against $N n_g+1$ rollouts for forward differences. The construction applies to A2 and A3 only. In A0 and A1 the algebraic variables are propagated by the tangent update and are not a function of $\delta$, and these variants keep finite-difference gradients in this study. The implementation is in \texttt{code/network\_gradient.py}; Section~\ref{sec:grad_results} compares it with central differences.

\subsection{Residual trust gate and dispatch reconciliation}
\label{sec:gate}
Section~\ref{sec:regime} shows that the learned correction can do harm when the dispatch changes without notice. The gate turns this observation into a rule, and a second step tries to repair the cause. At sample $k$ the gate computes the one-step innovation $s_k$ of the full corrected predictor: the RMS difference between the predicted and the measured next state (angles relative to the centre of inertia, and speeds), starting from the measured state at $k-1$ with the applied input and the known load. The score $q_k$ is the mean of the last three values of $s$ divided by $q_0$, the split-conformal $95\%$ quantile, over calibration trajectories, of the largest smoothed innovation in ordinary operation (measurement noise $10^{-3}$ on every state). The trust and the effective residual are
\begin{equation}
\begin{split}\tau_k&=\exp\{-a\max(q_k-1,\,0)\}\in(0,1],\\ d_{\rm eff}&=\tau_k\,\hat d .\end{split}
\label{eq:gate}
\end{equation}
The dead zone at $q_k=1$ keeps $\tau=1$ in ordinary operation, and $a$ is chosen on validation scenarios that are not used for testing. A candidate second score, based on the regularity margin $\eta_J$, was tried in the same selection; Section~\ref{sec:gate_results} reports that it did not help, so the operational gate uses the innovation only. The Newton projection stays active for every $\tau$, and $\tau=0$ gives the physics model with projection.

\emph{Reconciliation.} The load is measured, so an unannounced change of dispatch is the main source of a persistent innovation. The predictor therefore carries a dispatch scale $\rho$, with $P_m^{\rm pred}=\rho P_m^{\rm nom}$ and $\rho=1$ at the start. When $q_k>1$ the scale is re-estimated by one Gauss--Newton step on the flagged innovations of the last three samples,
\begin{equation}
\rho\leftarrow\rho-\frac{\sum_j g_j^\top e_j}{\sum_j g_j^\top g_j},\qquad g_j=\frac{e_j(\rho+\delta\rho)-e_j(\rho)}{\delta\rho},
\label{eq:rho}
\end{equation}
where $e_j(\rho)$ is the innovation vector of interval $j$ with the dispatch scale $\rho$. The predictor then uses the new dispatch, the innovation shrinks, $q$ falls below one when the three-sample window has cleared, and the trust returns. If the dispatch later returns to its old value the same mechanism moves $\rho$ back. The estimator assumes that the dispatch moves as a common scale, as it does when it follows the load; it knows nothing about individual generators, and it absorbs any other persistent bias in the same way. The reconciliation should therefore be read as a consistency-restoration step for the declared common-scale dispatch uncertainty, not as a general fault or parameter estimator. Both the gate and the reconciliation are heuristics without a guarantee.

\section{Analytical results}
\label{sec:theory}
The four propositions below are meant to be read as one chain, regular algebraic manifold, bounded manifold sensitivity, Lipschitz reduced dynamics, bounded growth of the residual error, stated in full after Proposition~\ref{prop:gronwall}. Each link has a computed counterpart (Table~\ref{tab:map}). None of them is a statement about the stability of the closed loop.

\begin{proposition}[Scalar benchmark is globally index one]
\label{prop:scalar}
For every $(\delta,u)\in\RR^2$ the equation $g(\delta,v,u)=0$ in \eqref{eq:s3} has exactly one real solution $v=h(\delta,u)$, $h$ is real-analytic, and $|\partial h/\partial\delta|\le a$.
\end{proposition}
\begin{proof}
For fixed $(\delta,u)$ the map $v\mapsto g$ has derivative $1+3c_bv^2\ge1$, so it is strictly increasing and unbounded in both directions, hence a bijection of $\RR$. The implicit function theorem gives analyticity, and $\partial h/\partial\delta=-a\cos\delta/(1+3c_bh^2)$ has modulus at most $a$.
\end{proof}
Because $c_b>0$, the elimination of $v$ is valid globally. The reduced model $\dot\delta=\omega$, $M\dot\omega=P-k_vh(\delta,u)\sin\delta-D\omega+u-c_3\delta^3+p$ is therefore defined for all states.

\begin{proposition}[Regularity margin of the network realization]
\label{prop:network}
Let the network graph be connected, let at least one generator be present, and let $B=\sum_l b_l a_la_l^\top+\sum_i b_i e_{\sigma(i)}e_{\sigma(i)}^\top$ with $a_l=e_j-e_k$ for the line $l=(j,k)$. Then $B\succ0$. Let $J=-\partial g/\partial\theta$ in \eqref{eq:n2}. If $|\theta_j-\theta_k|\le\gamma$ on every line and $|\delta_i-\theta_{\sigma(i)}|\le\gamma$ for every generator, with $\gamma<\pi/2$, then
\begin{equation}
\begin{split}&\lambda_{\min}(J)\ge\cos\gamma\,\lambda_{\min}(B)>0,\\ &\|J^{-1}\|_2\le \frac{1}{\cos\gamma\,\lambda_{\min}(B)} .\end{split}
\label{eq:cert}
\end{equation}
\end{proposition}
\begin{proof}
Differentiating \eqref{eq:n2} gives $J=\sum_l b_l\cos(\theta_j-\theta_k)\,a_la_l^\top+\sum_i b_i\cos(\delta_i-\theta_{\sigma(i)})\,e_{\sigma(i)}e_{\sigma(i)}^\top$. Each coefficient lies in $[b\cos\gamma,\,b]$ and each rank-one matrix is positive semidefinite, so $J\succeq\cos\gamma\,B$. For $B\succ0$ suppose $B\theta=0$ with $\theta\ne0$; then $\theta^\top B\theta=\sum_l b_l(\theta_j-\theta_k)^2+\sum_i b_i\theta_{\sigma(i)}^2=0$ forces equal angles on connected buses and a zero value at a generator bus, hence $\theta=0$.
\end{proof}
As a result the algebraic variables are locally unique and $C^1$ in $(\delta,P_d)$ on the region defined by the angle limits, and the Newton corrector has an invertible Jacobian there. The ratio $\eta_J=\lambda_{\min}(J)/(\cos\gamma\,\lambda_{\min}(B))$ is used below as a single regularity coordinate: $\eta_J>1$ means that the computed smallest eigenvalue lies above the topological bound. It describes the operating point and is not claimed to cause learning to fail. The bound is conservative (Table~\ref{tab:reg}) but cheap: $\lambda_{\min}(B)$ is computed once from the topology. The proposition says nothing about the stability of the closed loop.

\begin{proposition}[Lipschitz constant of the reduced network dynamics]
\label{prop:lip}
Under the assumptions of Proposition~\ref{prop:network}, let $E\in\RR^{n\times m}$ be the generator--bus incidence matrix ($E_{\sigma(i)i}=1$, other entries zero), $k$ the largest number of generators at one bus, $b_{\max}=\max_i b_i$ and
\begin{equation}
\kappa_\gamma=\frac{k\,b_{\max}}{\cos\gamma\,\lambda_{\min}(B)} .
\label{eq:kappa}
\end{equation}
In the region defined by the angle limits, $\theta=\theta(\delta)$ is $C^1$, $\|\partial\theta/\partial\delta\|_2\le\kappa_\gamma$, and the rotor-to-bus angles $\psi=\delta-E^\top\theta$ satisfy $\|\partial\psi/\partial\delta\|_2\le1+\kappa_\gamma$. It follows that the vector field of the reduced physics model, $\dot\delta=\omega$, $M\dot\omega=P_m+u-D\omega-(b_i\sin\psi_i)_i$, has Jacobian norm at most
\begin{equation}
L_f=1+\max_i\frac{D_i}{M_i}+\max_i\frac{b_i}{M_i}\,(1+\kappa_\gamma),
\label{eq:Lf}
\end{equation}
and is Lipschitz with this constant in $(\delta,\omega)$ on every convex subset of the region.
\end{proposition}
\begin{proof}
Differentiating $g(\delta,\theta(\delta))=0$ gives
\[J\,\partial\theta/\partial\delta=\partial g/\partial\delta=EC,\]
with $C=\mathrm{diag}(b_i\cos\psi_i)$, so $\partial\theta/\partial\delta=J^{-1}EC$. By Proposition~\ref{prop:network}, $\|J^{-1}\|_2\le1/(\cos\gamma\,\lambda_{\min}(B))$; moreover $\|E\|_2=\sqrt k$ and $\|C\|_2\le b_{\max}$, so that, because $k\ge1$,
\[\|\partial\theta/\partial\delta\|_2\le\frac{\sqrt k\,b_{\max}}{\cos\gamma\,\lambda_{\min}(B)}\le\kappa_\gamma .\]

Then $\partial\psi/\partial\delta=I-E^\top\partial\theta/\partial\delta$ and $\|E^\top\|_2=\sqrt k$ give
\[\|E^\top\partial\theta/\partial\delta\|_2\le\frac{k\,b_{\max}}{\cos\gamma\,\lambda_{\min}(B)}=\kappa_\gamma ,\]
which is the second bound.

The Jacobian of the reduced vector field has the block form $\bigl(\begin{smallmatrix}0&I\\A&-M^{-1}D\end{smallmatrix}\bigr)$ with $A=-M^{-1}\mathrm{diag}(b_i\cos\psi_i)\,\partial\psi/\partial\delta$, whose norm is at most $1+\|A\|_2+\|M^{-1}D\|_2$ and $\|A\|_2\le\max_i(b_i/M_i)(1+\kappa_\gamma)$.
\end{proof}
The sensitivity bound can be compared with the actual sensitivity. At $400$ states on the manifold for each case, drawn around the nominal equilibrium and kept within $\gamma=30^\circ$ on every line and rotor, the largest observed $\|\partial\theta/\partial\delta\|_2$ is \chainA{} (case9) and \chainB{} (case39), against $\kappa_\gamma=\chainC{}$ and $\chainD{}$; a finite-difference derivative of the Newton solution agrees with $J^{-1}EC$ to better than $10^{-6}$ (\texttt{tools/regularity\_chain\_check.py}). The bound is conservative by a factor of three to five, which is the price of obtaining it from the topology alone.

\begin{proposition}[Propagation of the residual error]
\label{prop:gronwall}
Let $F(x,t)$ be the vector field of the differential states of the plant and $\hat F(x',t)$ that of the predictor, in each case with the algebraic variables obtained from the respective constraint $g=0$: for the plant with the true exogenous input, for the predictor with the input it assumes. Suppose that along the two trajectories
\begin{equation}
\|F(x,t)-\hat F(x',t)\|\le L\|x-x'\|+\rho+\bar w ,
\label{eq:gron_hyp}
\end{equation}
where $\rho$ bounds the error of the learned residual and $\bar w$ the effect of the exogenous disturbance, both in units of the state derivative. Then the state error $e(t)=x(t)-x'(t)$ with $e(0)=0$ obeys $\|e(t)\|\le(\rho+\bar w)\,(e^{Lt}-1)/L$.
\end{proposition}
\begin{proof}
Gronwall's inequality applied to $\|e\|\le\int_0^t(L\|e\|+\rho+\bar w)\,\dd s$.
\end{proof}
Proposition~\ref{prop:gronwall} is an abstract propagation lemma, and what $L$, $\rho$ and $\bar w$ are depends on the benchmark. In the scalar benchmark the disturbance acts on the frequency equation only, so $\bar w=A/M$. In the network realization there is no additive disturbance. The predictor holds the load at its value at the start of the prediction window, and the load fluctuation reaches the dynamics through the algebraic law. By the implicit function theorem, $\partial\theta/\partial P_d=-J^{-1}$, so Proposition~\ref{prop:network} limits the change of the bus angles caused by a load deviation $\Delta P_d$ to $\|\Delta P_d\|_2/(\cos\gamma\,\lambda_{\min}(B))$ on a convex subset of the region. Since $\psi=\delta-E^\top\theta$, the acceleration changes by at most $\bar w=\sqrt{k}\,\max_i(b_i/M_i)\,\|\Delta P_d\|_2/(\cos\gamma\,\lambda_{\min}(B))$, and this value can be used in \eqref{eq:gron_hyp}. The estimate follows from Proposition~\ref{prop:network}; it is not a separate result, it is not evaluated in the experiments, and the empirical envelope does not use it.

The bound shows why the projection matters for this analysis. With exact algebraic variables the only error sources are $\rho$ and $\bar w$, whereas a tangent update adds a term of order $\|\Delta x\|^2$ per step that $\rho$ does not cover. Together with Proposition~\ref{prop:lip}, hypothesis \eqref{eq:gron_hyp} holds with $L=L_f+L_d$ for the network realization, where $L_d$ is a Lipschitz constant of the true residual and the segment between the two trajectories stays in the region of the angle limits. The regularity margin $\cos\gamma\,\lambda_{\min}(B)$ then enters the growth rate $e^{Lt}$ of the error bound. The chain is: regular algebraic manifold, bounded manifold sensitivity, Lipschitz reduced dynamics, bounded error growth. The resulting bound is far more conservative than the empirical envelope of Section~\ref{sec:tube}, which is why the controller uses the latter. No closed-loop Lyapunov certificate is claimed.

\section{Experimental protocol}
\label{sec:protocol}
All experiments are deterministic: every random draw uses a fixed seed (listed in the scripts), integration is explicit Runge--Kutta of order four with the step sizes stated in \ref{app:params}, and the full pipeline is run by one command (see the data and code availability statement). Prediction metrics are computed on windows of $10$ steps ($1$\,s) taken every $4$ (network) or $3$ (scalar) samples from closed-loop test runs under the static law, with the recorded inputs replayed to every predictor, so that the predictors are compared on identical data. The state RMSE combines angles (relative to the centre of inertia in the network case) and speeds. The algebraic residual is the actual constraint, $\max_j|g_j|$ or $|g|$ evaluated on the predicted state, and not a surrogate. Closed-loop runs use the true plant and a different set of initial conditions and disturbance phases for each realization; confidence intervals are percentile bootstrap intervals ($4000$ resamples) over realizations, and comparisons between controllers are paired. Counts of realizations are modest (Table~\ref{tab:scalar_cl}, Table~\ref{tab:net_cl}) and the intervals reflect that.

\section{Results}
\label{sec:results}
\subsection{Scalar benchmark: nonlinear structure}
The closed-loop equilibrium at the set point is $\delta^\ast=0.530$ with eigenvalues $-0.68\pm1.73\,i$. Continuing the equilibrium in a constant force bias (Fig.~\ref{fig:s_nl}a) shows that the feedback law moves the equilibrium less than the open loop and flattens the branch for negative bias; the kink near $p\approx0.13$ is the input saturation. Over a grid of initial conditions in $\delta_0\in[-0.6,2.6]$, $\omega_0\in[-2.2,2.2]$, 100\% of the trajectories converge to the nominal equilibrium with feedback and 29\% without it (Fig.~\ref{fig:s_nl}b,c). This is a numerical basin estimate on a finite grid and horizon ($20$\,s), not a proof.
\begin{figure*}[!t]\centering
\includegraphics[width=\linewidth]{fig_scalar_nonlinear.pdf}
\caption{Scalar benchmark. (a) Equilibria versus constant force bias with stability from the finite-difference Jacobian. (b,c) Convergence to the nominal equilibrium from a $64\times64$ grid of initial conditions, with and without feedback.}
\label{fig:s_nl}\end{figure*}

\subsection{Scalar benchmark: residual learning}
\label{sec:scalar_results}
The library has $20$ terms (degree $3$ selected on validation) and the parsimony rule keeps $3$. On the held-out test set the residual has RMS $0.120$ and the error after correction has RMS $0.021$, a reduction of 82\% (Fig.~\ref{fig:s_learn}). The validation criterion works on noisy weak-form targets and cannot tell the three-term model from much denser ones, but the exact held-out error can: a fit with $15$ to $20$ active terms reaches 0.0136, against 0.0212 for four terms (Supplementary Material S1). The parsimony rule therefore buys interpretability at the price of some accuracy, and all scalar results below use the three-term model. The one-step bound from a pointwise quantile is $\rho_{\rm pt}=0.110$ and covers a fraction 0.933 of independent test points, below the nominal $0.95$ because neighbouring samples are dependent; the block-maximum bound is $\rho_{\rm blk}=0.161$ with coverage 0.958 on 96 independent test blocks. The learning curve (\ref{app:extra}) shows that $2$, $8$ and $32$ training trajectories give a held-out error of 0.057, 0.020 and 0.011 at noise level $10^{-3}$, and 0.126, 0.031 and 0.015 at $10^{-2}$.
\begin{figure*}[!t]\centering
\includegraphics[width=\linewidth]{fig_scalar_learning.pdf}
\caption{Scalar benchmark, residual identification. (a) Learned versus true residual on held-out points. (b) Distribution of the measured residual error with the pointwise bound $\rho_{\rm pt}$. (c) Spatial error of the fit at $u=-0.5$.}
\label{fig:s_learn}\end{figure*}

\subsection{Scalar benchmark: prediction ablation}
Over the $1.2$\,s horizon the physics-only predictor has state RMSE 0.106, the corrected predictors 0.073 (A1 and A2 agree to four digits: 0.073 and 0.073). The per-realization ratio of A2 to A0 has mean 0.59 with bootstrap interval [0.47, 0.70] over $16$ test realizations. The projection does not change the accuracy here because the algebraic variable enters the acceleration only weakly; it changes the consistency. The mean worst-case constraint residual over the prediction window is $2.6\times10^{-5}$ for A0, $2.4\times10^{-5}$ for A1 and $1.1\times10^{-16}$ for A2, with maxima $2.6\times10^{-4}$, $2.3\times10^{-4}$ and $2.2\times10^{-16}$ (Fig.~\ref{fig:s_abl}). The residual of A2 is round-off. The consistency gain of the projection is therefore real but small in absolute terms in this benchmark.
\begin{figure*}[!t]\centering
\includegraphics[width=0.85\linewidth]{fig_scalar_ablation.pdf}
\caption{Scalar benchmark, prediction ablation on held-out closed-loop runs. (a) State RMSE versus prediction horizon for A0--A2. (b) Paired per-realization RMSE for A0 and A2 on a logarithmic axis.}
\label{fig:s_abl}\end{figure*}

\subsection{Scalar benchmark: closed loop}
Table~\ref{tab:scalar_cl} and Fig.~\ref{fig:s_cl} compare the controllers over $16$ realizations (disturbance amplitude $0.02$--$0.12$, random frequency and phase). The receding-horizon controllers reduce the integral absolute tracking error of the static law by 46\% (A0 against C0). The learned residual does not reduce it further: the paired difference A2$-$A0 is -0.0012\ [-0.0094,\,+0.0060], an interval that contains zero. The mean absolute error over the last three seconds improves slightly with the residual ($-0.0010\ [-0.0016,\,-0.0004]$), but the constraint-violation integral becomes larger ($+0.0052\ [+0.0028,\,+0.0077]$ for A2$-$A0). With the penalty tightening (A3) the violation falls below the A0 level ($-0.0071\ [-0.0120,\,-0.0026]$ for A3$-$A0, an interval that excludes zero) and the tracking error is not worse than for A0 ($-0.0071\ [-0.0181,\,+0.0067]$, interval containing zero). One possible explanation is that a more accurate predictor lets the optimizer approach the limit more closely and the tightening compensates; the experiments do not isolate that mechanism. The disturbance sweep in \ref{app:extra} shows that the advantage of A3 over A0 vanishes at larger disturbance amplitudes. The envelope used by A3 was calibrated on $20$ trajectories ($10$ static, $10$ MPC). On $16$ independent static-law trajectories it covered all of them at step $6$ and $0.875$ of them jointly over all twelve steps; on $16$ MPC trajectories the coverage was $1.000$ at every step. These values are not far from the nominal level, but with $16$ trajectories a fraction of $0.875$ is statistically compatible with $0.95$. The computing time per solve is 5.9\,ms for A2 against 2.0\,ms for A0.
\begin{table*}[!t]\centering\small
\caption{Scalar benchmark, closed loop, mean over $16$ realizations. IAE is the integral absolute error of $\delta-\delta_{ref}$; Viol.\ is the integral of the excess of $\delta$ over its limit; $t_{\rm solve}$ is the mean time per control step.}
\label{tab:scalar_cl}
\begin{tabular}{lcccccc}
\toprule
 & IAE & 95\% CI (IAE) & Last 3\,s err. & Viol.\ & Peak $\delta$ & $t_{\rm solve}$ [ms]\\
\midrule
C0 (static) & 0.746 & [0.673, 0.815] & 0.034 & 0.036 & 0.699 & 0.0\\
A0 & 0.403 & [0.364, 0.442] & 0.009 & 0.011 & 0.651 & 2.0\\
A1 & 0.401 & [0.362, 0.442] & 0.008 & 0.016 & 0.662 & 3.4\\
A2 & 0.401 & [0.362, 0.442] & 0.008 & 0.016 & 0.662 & 5.9\\
A3 & 0.395 & [0.356, 0.435] & 0.008 & 0.004 & 0.637 & 13.1\\
\bottomrule
\end{tabular}
\end{table*}
\begin{figure*}[!t]\centering
\includegraphics[width=\linewidth]{fig_scalar_closedloop.pdf}
\caption{Scalar benchmark, closed loop. (a) One realization; dotted line: set point, dashed: angle limit. (b) Applied input. (c) IAE over $16$ realizations.}
\label{fig:s_cl}\end{figure*}

\subsection{Network realizations: regularity and spectrum}
Table~\ref{tab:reg} lists the algebraic-Jacobian quantities at the nominal equilibrium. Continuing the equilibrium in a uniform load scale $\mu$ (Fig.~\ref{fig:n_reg}), the smallest eigenvalue of $J$ decreases monotonically while the maximum line angle grows, and the continuation ends at $\mu=4.500$ (case9) and $\mu=3.125$ (case39), the last scale at which the equilibrium solver converged. This is a numerical loadability estimate for the lossless model and should not be read as an operational limit. The linearization of the true plant at the equilibrium is stable for both cases, with slowest damping ratios $0.28$ and $0.24$ and mode frequencies between $1.36$ and $1.39$\,Hz (case9) and $1.15$ and $1.70$\,Hz (case39).
\begin{table*}[!t]\centering\small
\caption{Regularity of the network realizations at the nominal equilibrium. $\lambda_{\min}(B)$ is the topological constant of Proposition~\ref{prop:network}; the last row is its bound \eqref{eq:cert} for $\gamma=30^\circ$.}
\label{tab:reg}
\begin{tabular}{lcc}
\toprule
 & case9 & case39\\
\midrule
equilibrium residual $\max|g|,\ \max|P_m-P_e|$ & $1.1\times10^{-15}$ & $6.2\times10^{-15}$\\
$\lambda_{\min}(J)$ at equilibrium & 2.24 & 4.19\\
condition number of $J$ & 25.6 & 246.6\\
max line angle at equilibrium [deg] & 7.8 & 9.6\\
$\lambda_{\min}(B)$ & 2.25 & 4.30\\
bound $\cos30^\circ\,\lambda_{\min}(B)$ & 1.95 & 3.73\\
smallest $\lambda_{\min}(J)$ seen in closed loop & 2.13 & 3.44\\
margin $\eta_J$ at equilibrium & 1.15 & 1.12\\
\bottomrule
\end{tabular}
\end{table*}
In case9 the smallest eigenvalue observed in any closed-loop run stays above the $30^\circ$ bound. In case39 it falls below it, which by \eqref{eq:cert} implies that some rotor-to-bus angle exceeded $30^\circ$ (the line-angle limit was respected; the rotor-angle limit is $30^\circ+0.35$\,rad); the observed value corresponds to a maximum angle of at least about $37^\circ$. The Jacobian remained well conditioned throughout.
\begin{figure*}[!t]\centering
\includegraphics[width=\linewidth]{fig_network_regularity.pdf}
\caption{Network realizations. (a) Regularity margin along a uniform load-scale continuation. (b) Maximum line angle along the continuation; dashed: $30^\circ$. (c) Eigenvalues of the linearized plant at the nominal equilibrium (the neutral mode of the angle shift is omitted).}
\label{fig:n_reg}\end{figure*}

\subsection{Network realizations: residual learning, prediction and algebraic consistency}
\label{sec:net_pred}
The residual features are fitted per generator. At the nominal loading the held-out residual has RMS 2.40 (case9) and 5.49 (case39) in acceleration units; after correction the error has RMS 0.058 and 0.088, reductions of 97.6\% and 98.4\%. Of the available library terms, 24 of 39 (case9) and 118 of 130 (case39) remain active, so these are thresholded regression models and not sparse ones. Raising the threshold shows how much of the fit is needed: with 12 active terms in case9 the held-out error reduction is still 97.1\% (against 97.1\% with all terms), and with 50 terms in case39 it is 98.1\% (against 98.5\%) (Supplementary Material S1). A genuinely sparse network model would need a smaller, physically motivated library.

Table~\ref{tab:net_pred} summarizes the prediction ablation on held-out closed-loop runs under load fluctuations of up to $8\%$, twice the training maximum. Over $1$\,s the physics-only predictor has state RMSE 0.704 (case9) and 1.631 (case39); the corrected predictors reach 0.159 and 0.538. The per-run ratio A2/A0 has mean 0.17 [0.09, 0.25] for case9 and 0.24 [0.17, 0.31] for case39. In a Monte Carlo ensemble of 24 realizations with load scale, line scale, inertia and damping drawn jointly from bounded ranges ($0.8$--$1.2$ for load, inertia and damping, $0.85$--$1.15$ for lines; the plant changes, the predictors keep their nominal parameters) the mean ratio is 0.16 [0.13, 0.20] and 0.16 [0.13, 0.19]. These intervals describe simulation variability under the declared distributions, not statistical inference about a physical system. Without the Newton corrector the algebraic residual grows along the prediction window (mean over windows of the largest $|g_j|$: $1.7\times10^{-3}$ for case9 and $2.3\times10^{-2}$ for case39, maxima $9.8\times10^{-2}$ and $6.4\times10^{-1}$); with it the residual is at the solver tolerance ($4.1\times10^{-11}$ and $4.1\times10^{-11}$). In contrast to the scalar benchmark, the projection also improves accuracy at short horizons (case9, $0.1$\,s: A1 $0.017$, A2 $0.011$; case39: $0.039$ and $0.021$), but A1 and A2 are close at $1$\,s.
\begin{table*}[!t]\centering\small
\caption{Network prediction ablation over a $1$\,s horizon (nominal loading, held-out runs). State RMSE, and the mean over windows of the largest algebraic residual.}
\label{tab:net_pred}
\begin{tabular}{llccc}
\toprule
Case & Model & State RMSE ($0.5$\,s) & State RMSE ($1$\,s) & $\max_j|g_j|$\\
\midrule
case9 & A0 & 0.592 & 0.704 & $1.7\times10^{-3}$\\
case9 & A1 & 0.127 & 0.151 & $1.7\times10^{-3}$\\
case9 & A2 & 0.131 & 0.159 & $4.1\times10^{-11}$\\
case39 & A0 & 1.297 & 1.631 & $2.3\times10^{-2}$\\
case39 & A1 & 0.296 & 0.525 & $2.3\times10^{-2}$\\
case39 & A2 & 0.297 & 0.538 & $4.1\times10^{-11}$\\
\bottomrule
\end{tabular}
\end{table*}
\begin{figure*}[!t]\centering
\includegraphics[width=\linewidth]{fig_network_ablation.pdf}
\caption{Network realizations, prediction. (a,b) State RMSE versus horizon for A0--A2. (c) Ratio of the A2 to the A0 RMSE in the Monte Carlo ensemble with bootstrap $95\%$ intervals.}
\label{fig:n_abl}\end{figure*}

\subsection{Network realizations: sensitivity and operating-point shift}
\label{sec:sens}
Two different tests are reported, and they answer different questions. In the first, one plant parameter is varied at a time while the predictors keep every nominal parameter, including the nominal dispatch, and the residual identified at the nominal loading (Table~\ref{tab:sens}). This test uses its own set of 5 closed-loop runs per scale (load fluctuation up to $6\%$, analysis windows every 6 samples), so its entries at scale 1.0 (0.16 and 0.30) differ slightly from the per-run means of Section~\ref{sec:net_pred}; the two protocols use different test runs. Scaling the line reactances, inertia and damping leaves the benefit intact (A2/A0 ratios between $0.17$ and $0.28$ for case9 and between $0.30$ and $0.38$ for case39). Scaling the load is different, because the plant dispatch follows the load while the predictor still assumes the nominal dispatch, which is an unmodelled change in mechanical power. At $0.6$ and $0.8$ the corrected predictor is still better than the physics-only one but the margin shrinks, and at $1.25$ and $1.5$ it is about twice as inaccurate. So the learned residual does not protect against, and can aggravate, an unmodelled shift of the power balance.

A natural explanation is that the residual is specific to the operating point at which it was fitted. A second test argues against it. Here the predictor is told the loading and the dispatch, and only the residual differs: one residual is fitted at the nominal loading and one is fitted on data pooled over five loadings ($0.6$ to $1.5$ times nominal). Both are evaluated at loadings seen in training, at interpolated loadings and at $1.8$ times nominal (Table~\ref{tab:transfer}). The A2/A0 ratio of the single-loading residual lies between 0.08 and 0.32 everywhere, including at $1.8$ times nominal, and pooled training changes it by less than 0.003 in absolute terms. In these reduced models the learned residual is therefore close to load-independent, because it is a function of the rotor angle, speed and input, and the loading enters through the known power balance. Whether this carries over to models with load dynamics or voltage-dependent loads is not tested here. Complete tables are in Supplementary Material S1.
\begin{table*}[!t]\centering\small
\caption{Ratio of the A2 to the A0 state RMSE ($1$\,s horizon) when the plant parameter is scaled and the predictors keep their nominal parameters, including the dispatch.}
\label{tab:sens}
\begin{tabular}{lccccc}
\toprule
 & \multicolumn{5}{c}{scale}\\
Case, parameter & 0.6 & 0.8 & 1.0 & 1.25 & 1.5\\
\midrule
case9, load & 0.60 & 0.45 & 0.16 & 1.96 & 1.70\\
case9, line reactance & 0.17 & 0.17 & 0.16 & 0.17 & 0.17\\
case9, inertia & 0.28 & 0.18 & 0.16 & 0.21 & 0.27\\
case9, damping & 0.28 & 0.18 & 0.16 & 0.21 & 0.27\\
\midrule
case39, load & 0.71 & 0.52 & 0.30 & 1.99 & 1.84\\
case39, line reactance & 0.30 & 0.30 & 0.30 & 0.30 & 0.30\\
case39, inertia & 0.38 & 0.31 & 0.30 & 0.32 & 0.36\\
case39, damping & 0.37 & 0.31 & 0.30 & 0.32 & 0.36\\
\bottomrule
\end{tabular}
\end{table*}
\begin{table*}[!t]\centering\small
\caption{Ratio of the A2 to the A0 state RMSE ($1$\,s horizon) when the predictor knows the loading and the dispatch. ``Single'': residual fitted at nominal loading; ``pooled'': residual fitted on five loadings. Loadings marked $^{\ast}$ were used for the pooled fit.}
\label{tab:transfer}
\begin{tabular}{lcccc}
\toprule
 & \multicolumn{2}{c}{case9} & \multicolumn{2}{c}{case39}\\
Load scale & single & pooled & single & pooled\\
\midrule
0.6$^{\ast}$ & 0.29 & 0.30 & 0.17 & 0.17\\
1.0$^{\ast}$ & 0.16 & 0.17 & 0.24 & 0.24\\
1.5$^{\ast}$ & 0.32 & 0.32 & 0.17 & 0.17\\
0.7 & 0.24 & 0.25 & 0.20 & 0.20\\
0.9 & 0.15 & 0.16 & 0.21 & 0.21\\
1.1 & 0.08 & 0.08 & 0.21 & 0.21\\
1.4 & 0.15 & 0.15 & 0.24 & 0.24\\
1.8 & 0.22 & 0.22 & 0.31 & 0.31\\
\bottomrule
\end{tabular}
\end{table*}

\subsection{Network realizations: a mismatch the library cannot represent exactly}
\label{sec:ool}
The saliency and friction terms of the network plant have the same form as library functions, which makes identification easier than it would be with unknown dynamics. To probe this, the plant was given one more term, $0.15\,b_i\psi_i|\psi_i|$ in the electrical power, which no combination of the library functions reproduces exactly. The residual was refitted with the same library and the same procedure. The held-out error reduction was 97.4\% (case9) and 99.0\% (case39), and the one-second prediction ratio A2/A0 was 0.17 [0.10, 0.25] and 0.22 [0.16, 0.28]. The library is flexible enough to approximate the extra term over the range of angles visited, so this is evidence of approximation capacity and not of recovery of the true structure; a discontinuous or strongly state-dependent effect would be a harder test.

\subsection{Network realizations: closed loop}
Table~\ref{tab:net_cl} compares the controllers at an elevated loading ($\mu=2.75$ for case9, $\mu=2.25$ for case39), where the equilibrium line angles are about $22^\circ$ and the residual is re-identified at that loading. The load fluctuates with amplitude $7$--$11\%$ at $2.5$--$4$\,rad/s. The receding-horizon controllers with the learned residual reduce the RMS speed deviation of the static damping-injection law by 15\% (case9) and 4\% (case39), whereas the physics-only controller A0 gives no reduction (case9) or a slightly worse value (case39). The peak line angle is not reduced by any controller: it is within a fraction of a degree of the static law, far from the $30^\circ$ limit. The input energy of the receding-horizon controllers is higher. A3 coincides with A2 to the printed precision, that is, the penalty tightening is inactive in these runs.
\begin{table*}[!t]\centering\small
\caption{Network closed loop at elevated loading; means over 6 (case9) and 4 (case39) realizations. Peak: largest line-angle difference; $\omega_{\rm rms}$: RMS speed deviation after the first third of the run; $E_u$: input energy normalized by the input bound.}
\label{tab:net_cl}
\begin{tabular}{llcccc}
\toprule
Case & Ctrl. & Peak [deg] & $\omega_{\rm rms}$ [rad/s] & $E_u$ & $t_{\rm solve}$ [ms]\\
\midrule
case9 & C0 & 23.02 & 1.026 & 10.2 & 0\\
 & A0 & 23.37 & 1.027 & 12.3 & 11\\
 & A2 & 23.08 & 0.876 & 16.0 & 16\\
 & A3 & 23.08 & 0.876 & 16.0 & 16\\
\midrule
case39 & C0 & 23.65 & 1.960 & 34.8 & 0\\
 & A0 & 23.68 & 2.067 & 43.2 & 123\\
 & A2 & 23.66 & 1.878 & 45.8 & 270\\
 & A3 & 23.66 & 1.878 & 45.8 & 270\\
\bottomrule
\end{tabular}
\end{table*}
To test the tightening where the limit is active (envelope from 7 trajectories, half of them under MPC), the case9 experiment was repeated with a line-angle limit of 23.0$^\circ$, just above the equilibrium angle, and the largest calibrated envelope (3.1$^\circ$ at the fifth step). Over $6$ realizations the peak angle was 23.02$^\circ$ (C0), 23.08$^\circ$ (A2) and 23.09$^\circ$ (A3), with the fraction of time above the limit 2.7\%, 3.6\% and 3.6\%. The tightening again changed nothing: with the input bound used here the generators cannot move the line angles materially within the short horizon, so the constraint is determined by the load. The network results therefore support the prediction and consistency claims and the speed-regulation gain, not a constraint-handling claim. On the envelope itself, calibration used 7 trajectories per case at the elevated loading (static law and MPC, same disturbance range as the closed-loop runs). On independent trajectories the step-5 envelope covered 1.00 (case9) and 1.00 (case39) of the static-law trajectories and 1.00 and 1.00 of the MPC trajectories, and 1.00 and 1.00 of the closed-loop realizations of Table~\ref{tab:net_cl}; with so few trajectories these fractions are coarse. With finite-difference gradients the computing time per control step is about 16\,ms for case9, below the sampling period of $100$\,ms, and about 270\,ms for case39, above it. Section~\ref{sec:grad_results} shows that the exact gradient of Section~\ref{sec:grad} removes this limit without changing the closed-loop result.
\begin{figure*}[!t]\centering
\includegraphics[width=\linewidth]{fig_network_closedloop.pdf}
\caption{Network realizations, closed loop at elevated loading: (a) case9, (b) case39. Top: maximum line angle, with the $30^\circ$ limit dashed. Bottom: maximum speed deviation. One realization per case; A2 and A3 coincide.}
\label{fig:n_cl}\end{figure*}

\subsection{Where the correction helps, where it stops, and what it costs}
\label{sec:regime}
The sensitivity, transfer and mismatch experiments together give a validity map for the correction (Table~\ref{tab:validity}). The boundary is not set by how far the loading lies from the one at which the residual was fitted. A residual fitted at nominal loading alone keeps the ratio between 0.08 and 0.32 up to $1.8$ times nominal when the predictor is told the loading and the dispatch (Table~\ref{tab:transfer}), whereas a predictor that keeps the nominal dispatch while the plant moves to $1.25$ or $1.5$ times nominal is about twice as inaccurate as the physics-only one. What decides the outcome is whether the known part of the model, here the power balance, is still right. The loading information is carried by the known power balance and not by the residual, and a residual cannot compensate for an error in the power balance. Between the two extremes the benefit shrinks gradually on the low-load side (ratio $0.45$ to $0.71$ at $0.8$ and $0.6$ times nominal load), while on the high-load side it is already lost at $1.25$; the crossing point was not located. Only one smooth extra term was tested as unmodelled structure (Section~\ref{sec:ool}); discontinuous or hidden-state effects were not, and no claim is made outside the range tested.

\begin{table*}[!t]\centering\small
\caption{Validity map of the learned correction, from Tables~\ref{tab:sens} and~\ref{tab:transfer} and Section~\ref{sec:ool}. Ratio of the A2 to the A0 state RMSE at $1$\,s, ranges over case9 and case39.}
\label{tab:validity}
\begin{tabular}{>{\raggedright\arraybackslash}p{4.6cm}>{\raggedright\arraybackslash}p{2.8cm}cl}
\toprule
Change of operating condition & Predictor is told & Ratio & Reading\\
\midrule
Line reactance, inertia or damping scaled by $0.6$--$1.5$ & nominal parameters & 0.16--0.38 & benefit intact\\
Loading $0.6$--$1.8$ times nominal & loading and dispatch & 0.08--0.32 & benefit intact\\
Smooth extra term outside the feature library & --- & 0.17, 0.22 & benefit intact\\
Load $0.6$ or $0.8$ times nominal, dispatch unannounced & nominal dispatch & 0.45--0.71 & benefit shrinks\\
Load $1.25$ or $1.5$ times nominal, dispatch unannounced & nominal dispatch & 1.70--1.99 & benefit reversed\\
\bottomrule
\end{tabular}
\end{table*}

The more faithful predictors also cost more, and the margin adds protection only where it is active. Table~\ref{tab:tradeoff} sets fidelity, protection and computing cost side by side. Going from A0 to A1 buys most of the accuracy. Going from A1 to A2 buys constraint consistency (the largest algebraic residual falls from $1.7\times10^{-3}$ and $2.3\times10^{-2}$ to $4.1\times10^{-11}$ on case9 and case39) at a higher computing cost. A3 adds the margin, which matters only on the scalar benchmark, where a step takes $13.1$\,ms against $5.9$\,ms for A2. On case39 the projected predictor with finite-difference gradients needs $270$\,ms per step against a sampling period of $100$\,ms. The exact gradient of Section~\ref{sec:grad}, for which the manifold sensitivity $J^{-1}EC$ of Proposition~\ref{prop:lip} is the main ingredient, lowers this to about $24$\,ms in a later timing session (Table~\ref{tab:grad}); the last row of Table~\ref{tab:tradeoff} gives that figure. A reduced-order solver has not been tried.

\begin{table*}[!t]\centering\small
\setlength{\tabcolsep}{5pt}
\caption{Fidelity and computing cost of the variants. State RMSE over $1.2$\,s (scalar) and $1$\,s (case9, case39) on held-out runs, and mean time per control step. A3 uses the A2 predictor. Dashes: not run in closed loop. $^{\dagger}$Timed in a later session on a busier machine (Table~\ref{tab:grad}); compare times within one session.}
\label{tab:tradeoff}
\begin{tabular}{lcccccc}
\toprule
 & \multicolumn{3}{c}{State RMSE} & \multicolumn{3}{c}{Time per step [ms]}\\
\cmidrule(lr){2-4}\cmidrule(lr){5-7}
Variant & scalar & case9 & case39 & scalar & case9 & case39\\
\midrule
A0 & 0.106 & 0.704 & 1.631 & 2.0 & 11 & 123\\
A1 & 0.073 & 0.151 & 0.525 & 3.4 & -- & --\\
A2 & 0.073 & 0.159 & 0.538 & 5.9 & 16 & 270\\
A3 & 0.073 & 0.159 & 0.538 & 13.1 & 16 & 270\\
A2, exact gradient$^{\dagger}$ & 0.073 & 0.159 & 0.538 & -- & 4.6 & 24\\
\bottomrule
\end{tabular}
\end{table*}

\subsection{Exact gradient and computing time}
\label{sec:grad_results}
The exact gradient of Section~\ref{sec:grad} was compared with central differences at random states and inputs of both networks, with and without active line-angle penalties. The largest relative difference is $1.2\times10^{-9}$ (case9) and $7.1\times10^{-9}$ (case39). One evaluation of cost and gradient takes $0.59$\,ms against $4.7$\,ms with forward differences on case9, and $4.6$\,ms against $149$\,ms on case39 (Table~\ref{tab:grad}).

Replacing the finite differences in A2 and A3 leaves the closed-loop result unchanged. On the realizations of Table~\ref{tab:net_cl} the RMS speed deviation of A2 is $0.8758$ (case9) and $1.8784$ (case39) with either gradient, and peak angle and input energy agree to at least five digits. The time per control step falls from $29$ to $4.6$\,ms on case9 and from $588$ to $24$\,ms on case39, and the largest step time observed on case39 is $28.5$\,ms, below the $100$\,ms sampling period. A3 behaves like A2, since its tightening is inactive. These runs were made on a single processor core in a later session, on a busier machine than the one used for Table~\ref{tab:net_cl}, so the absolute finite-difference times differ from those in the earlier table ($270$\,ms for A2 on case39). The ratio within one session is the quantity to compare, and absolute times will change with the hardware. The comparison concerns computing time only. A0 still uses finite differences, so its time does not show what an optimized physics-only controller would need, and the closed-loop gain of the projected predictor is as modest as before.

\begin{table*}[!t]\centering\small
\setlength{\tabcolsep}{4pt}
\caption{Exact gradient of the projected predictor against finite differences. Error: largest relative difference to central differences. Evaluation: time of one cost-and-gradient call. Step: mean time per control step of A2 in closed loop (6 realizations for case9, 4 for case39). $\omega_{\rm rms}$: RMS speed deviation of A2.}
\label{tab:grad}
\begin{tabular}{lcccccccc}
\toprule
 & & \multicolumn{2}{c}{Evaluation [ms]} & \multicolumn{2}{c}{Step [ms]} & \multicolumn{2}{c}{$\omega_{\rm rms}$ [rad/s]}\\
\cmidrule(lr){3-4}\cmidrule(lr){5-6}\cmidrule(lr){7-8}
Case & Error & exact & finite diff. & exact & finite diff. & exact & finite diff.\\
\midrule
case9 & $1.2\times10^{-9}$ & 0.59 & 4.7 & 4.6 & 28.6 & 0.8758 & 0.8758\\
case39 & $7.1\times10^{-9}$ & 4.6 & 149 & 24.0 & 588 & 1.8784 & 1.8784\\
\bottomrule
\end{tabular}
\end{table*}

\subsection{Validity atlas and gate boundary}
\label{sec:atlas}
The two boundaries found so far, the operating-condition boundary of Tables~\ref{tab:sens} and~\ref{tab:transfer} and the regularity boundary of Fig.~\ref{fig:n_reg}, are combined in Fig.~\ref{fig:atlas}. The horizontal axis is the load scale $\mu$ of the plant, labelled also with the regularity margin $\eta_J$; the vertical axis is the dispatch error $\kappa$, with the predictor told the load and using the dispatch $P_m(1+\kappa)$, so that $\kappa=0$ is a known dispatch. Each cell gives the ratio of the A2 to the A0 error over four runs, with the residual fitted at nominal loading. With a known dispatch the benefit survives the whole range up to $\mu=2.5$ (ratios 0.13--0.28 for case9 and 0.18--0.32 for case39), and it weakens only at the edge of the feasible range: on case39 at $\mu=3.0$, where $\eta_J=0.73$ and the largest line angle is $29.9^\circ$, the ratio is $0.82$. The margin $\eta_J$ falls below one on case39 already at $\mu=2.5$ ($0.91$), where the benefit is intact, so $\eta_J$ is a description of the operating point and not a predictor of failure. The dispatch axis matters more. A predictor that is $10\%$ short of the dispatch is worse than the physics model in every cell (ratios 1.78--2.96 for case9 and 1.65--3.65 for case39 up to $\mu=2.5$). An excess of $10\%$ or $20\%$ shrinks the benefit but does not remove it, and an excess of $40\%$ takes it away at the highest loadings. I did not isolate why a shortfall is so much worse than an excess. The atlas describes the reduced models only, and each cell rests on four runs.

The black corners in Fig.~\ref{fig:atlas} mark the cells in which the gate of Section~\ref{sec:gate}, run on the steady mismatch of that cell, falls back in more than half of the samples. It flags 17 of the 18 cells in which A2 is worse than A0 on case9 and all 19 on case39; the one miss is $\mu=0.6$ with $\kappa=-0.1$, where the absolute error of the dispatch is small. It does not flag any of the six cells with a known dispatch on case9, and flags one on case39 ($\mu=3.0$). The gate also flags most cells in which A2 is still better than A0 (17 of 24 on case9, 18 of 23 on case39), namely those with a dispatch excess. It detects a mismatch of the power balance, not the harm it does: the gate is a monitor of the consistency of the assumed operating condition and not a predictor of future error, which is why it errs on the cautious side.

\begin{figure*}[!t]\centering
\includegraphics[width=\linewidth]{fig_validity_atlas.pdf}
\caption{Validity atlas. Ratio of the A2 to the A0 one-second prediction error (numbers in the cells, colour: $\log_2$ of the ratio; blue: the correction helps, red: it hurts) against load scale $\mu$ (bracket: regularity margin $\eta_J$) and dispatch error $\kappa$ of the predictor. Black corner: the trust gate falls back in more than half of the samples.}
\label{fig:atlas}
\end{figure*}

\subsection{Trust gate, reconciliation and recovery}
\label{sec:gate_results}
The gate was tested on runs in which the load and the dispatch of the plant step together at $3$\,s, to a scale $\mu$ between $0.8$ and $1.5$ and with the predictor unaware of the dispatch, and on runs in which the step is undone at $6$\,s. The constant chosen on the validation scenarios (announced and unannounced steps, three runs each) is $a=0.25$. The regularity weight that was tried alongside it came out as zero, and the validation score is flat in it (best and worst setting differ by 0.3--0.4\%): a high loading with a known dispatch also has a low margin and the correction still works there, so the margin cannot tell the good cases from the bad ones. The calibrated levels are $q_0=0.024$ (case9) and $0.037$ (case39). Table~\ref{tab:rec} compares five predictors: A0, A2, A2 with the true dispatch (oracle), the gate, and the gate with reconciliation.

In ordinary operation, with stronger load fluctuation and with noise three times larger, the gate never fell below $\tau=0.5$ and $\rho$ stayed within $0.005$ of one, so the gated predictors equal A2 up to a few percent; the one visible cost is on case39 with the larger noise, where the reconciled error is 6\% above that of A2 (0.458 against 0.431). After an unannounced step all shifts were detected, within $0.1$ to $0.3$\,s (a $3\%$ step took $0.23$ to $0.28$\,s, the others $0.1$ to $0.18$\,s). In a further test the plant had errors of $15$ to $20\%$ in line reactance, inertia or damping that the predictor and the gate did not know. About $1\%$ of the samples at most fell below $\tau=0.5$, the mean trust stayed between $0.86$ and $1.00$ and $\rho$ moved by at most $0.08$ (Supplementary Material S1), so the gate responds to a persistent mismatch of the power balance and only mildly to these parametric errors; a line-reactance error of $15\%$ lowers the trust to about $0.86$ to $0.90$ without a fallback. The plain gate lowers the trust and returns the error to the physics-only level, which removes the error that A2 alone adds when A2 is worse than A0 (steps of $10\%$ and more upward) but gives up a benefit that A2 still has after the step down to $0.8$. With reconciliation the estimate $\rho$ follows the true scale to within $0.5\%$ (for example $1.251$ for a true $1.25$ and $1.504$ for $1.5$), the trust is back at $0.8$ within $0.5$\,s, and the error is close to that of the oracle (case9: 0.108 against 0.101 for $\mu=1.25$, 0.111 against 0.107 for $1.5$; case39: 0.479 against 0.453 and 0.321 against 0.310). When the step is undone, the dispatch scale goes back to one, the trust dips again for a few samples and returns, and the error after the recovery equals that of A2 (case9: 0.097 against 0.098 for a step of $10\%$ and 0.140 against 0.136 for $25\%$; case39: 0.152 against 0.150 and 0.236 against 0.222). The correction is therefore not locked out by a past shift. Fig.~\ref{fig:recovery} shows one run of the $25\%$ case.

A held-out family tests the common-scale assumption itself. In it the load steps by a common factor, but after the step each generator's dispatch is the scaled dispatch times a random factor between $0.8$ and $1.2$, renormalized to keep the power balance (six runs for case9, four for case39; Supplementary Material S1). All steps were detected within $0.1$\,s and $\hat\rho$ followed the common factor ($1.09$ and $1.25$ for steps of $10\%$ and $25\%$ on case9). On case9 the reconciled predictor stays well below A2 (0.222 against 0.727 for the $25\%$ step) but above the oracle (0.144), and the plain gate is better for the $10\%$ step (0.129 against 0.187). On case39 the reconciliation does worse than the plain gate (0.796 against 0.205 for the $10\%$ step, and 0.940 against 0.861 for $25\%$): the one-step innovation shrinks once a common scale is fitted, the trust returns, and the remaining non-uniform error is carried through the one-second prediction. The gate detects such a shift, but the reconciliation is a repair only for the declared common-scale case. When the structure is violated the plain fallback is the safer response, and the experiment is a reason to read the reconciliation narrowly.

\begin{table*}[!t]\centering\small
\setlength{\tabcolsep}{3.5pt}
\caption{Unannounced steps of load and dispatch at $3$\,s: one-second prediction RMSE from $0.6$\,s after the step while the step lasts, for A0, A2, A2 with the true dispatch (oracle), the trust gate and the gate with reconciliation; estimated dispatch scale $\hat\rho$; time from the step until $q>1$ (all steps were detected); lowest trust of the reconciling gate. Six runs per scenario for case9, four for case39.}
\label{tab:rec}
\begin{tabular}{llccccccc}
\toprule
Case & Step $\mu$ & A0 & A2 & Oracle & Gate & Reconc. & $\hat\rho$ & Delay [s] / $\tau_{\min}$\\
\midrule
case9 & 1.05 & 0.350 & 0.175 & 0.116 & 0.143 & 0.127 & 1.050 & 0.18 / 0.89\\
 & 1.10 & 0.221 & 0.300 & 0.131 & 0.169 & 0.141 & 1.104 & 0.10 / 0.59\\
 & 0.8 & 1.166 & 0.543 & 0.134 & 1.046 & 0.138 & 0.799 & 0.10 / 0.30\\
 & 1.25 & 0.370 & 0.725 & 0.101 & 0.407 & 0.108 & 1.251 & 0.10 / 0.20\\
 & 1.5 & 1.200 & 1.992 & 0.107 & 1.203 & 0.111 & 1.499 & 0.10 / 0.03\\
case39 & 1.05 & 0.800 & 0.440 & 0.316 & 0.353 & 0.317 & 1.049 & 0.10 / 0.73\\
 & 1.10 & 0.436 & 0.713 & 0.268 & 0.285 & 0.269 & 1.101 & 0.10 / 0.43\\
 & 0.8 & 2.568 & 1.293 & 0.143 & 2.495 & 0.144 & 0.800 & 0.10 / 0.15\\
 & 1.25 & 0.992 & 1.893 & 0.453 & 1.011 & 0.479 & 1.251 & 0.10 / 0.09\\
 & 1.5 & 2.670 & 4.998 & 0.310 & 2.671 & 0.321 & 1.504 & 0.10 / 0.01\\
\bottomrule
\end{tabular}
\end{table*}

\begin{figure*}[!t]\centering
\includegraphics[width=\linewidth]{fig_recovery.pdf}
\caption{Recovery test, step to $\mu=1.25$ at $3$\,s undone at $6$\,s (shaded). Top: true and estimated dispatch scale. Middle: trust of the plain gate and of the gate with reconciliation. Bottom: one-second prediction RMSE of A0, A2, A2 with the true dispatch, the gate and the reconciling gate.}
\label{fig:recovery}
\end{figure*}

\subsection{Mild unannounced step in closed loop}
\label{sec:gate_cl}
A large step drives the input to its bound in every predictive controller, and in the earlier tests with a step to $\mu=1.5$ all of them gave the same speed deviation (Supplementary Material S1). A mild step leaves the controller some freedom. Table~\ref{tab:gate_cl} gives the closed-loop result for steps of $5\%$ and $10\%$ in load and dispatch, with the controllers of Section~\ref{sec:grad} (exact gradients), an angle reference that follows the new equilibrium, and the dispatch unknown to the controller; the gate sees noisy measurements and the controller exact states. On case9 the effect is clear. With the wrong dispatch A2 is, for the $10\%$ step, worse than the physics-only A0 (speed deviation 0.140 against 0.114). The gate and the gate with reconciliation both bring it to the level of the oracle controller that knows the dispatch (0.011 to 0.023 against 0.015 to 0.017). The fallback state alone (physics model with projection, $\tau=0$) gives the same speed deviation as A0 (0.178 and 0.114), so the gain comes from keeping the correction when it is valid and withdrawing it only when it is not. The plain gate was as good as the reconciling one here, and better in the $10\%$ run, so the reconciliation improves the prediction (Table~\ref{tab:rec}) but did not improve the closed loop in these runs. A closed-loop recovery run on case9 (step of $10\%$ at $3$\,s, undone at $6$\,s, six realizations) shows the same ordering and the return of the correction: the speed deviation during the step is 0.139 for A2, 0.112 for the fallback alone, 0.012 for the gate, 0.023 for the reconciling gate and 0.017 for the oracle, and after the step is undone it is 0.017, 0.237, 0.015, 0.021 and 0.015, with the mean trust of both gates at 0.99 or higher. On case39 the input is at its bound for 60 to 100\% of the time in every predictive controller, the speed deviations of A2, the gates and the oracle differ by less than 2\%, and the closed loop cannot separate them, as in the large step.

\begin{table*}[!t]\centering\small
\setlength{\tabcolsep}{4pt}
\caption{Closed loop after an unannounced step at $3$\,s: RMS speed deviation [rad/s] from $0.5$\,s after the step, and share of samples with the input at its bound. Six realizations for case9, four for case39. Physics + proj.: the fallback state alone ($\tau=0$), run on case9 only. Oracle: A2 told the true dispatch.}
\label{tab:gate_cl}
\begin{tabular}{llcccccc}
\toprule
 & & \multicolumn{2}{c}{case9} & \multicolumn{2}{c}{case39}\\
\cmidrule(lr){3-4}\cmidrule(lr){5-6}
Step $\mu$ & Controller & $\omega_{\rm rms}$ & bound & $\omega_{\rm rms}$ & bound\\
\midrule
1.05 & C0 & 0.344 & 0.00 & 0.767 & 0.00\\
 & A0 & 0.178 & 0.22 & 0.541 & 0.60\\
 & A2 & 0.073 & 0.19 & 0.458 & 1.00\\
 & physics + proj. & 0.178 & 0.22 & -- & --\\
 & gate & 0.017 & 0.08 & 0.463 & 0.81\\
 & gate + reconc. & 0.018 & 0.03 & 0.462 & 0.86\\
 & oracle & 0.015 & 0.03 & 0.462 & 0.85\\
1.10 & C0 & 0.361 & 0.00 & 0.802 & 0.00\\
 & A0 & 0.114 & 0.17 & 0.544 & 0.74\\
 & A2 & 0.140 & 0.43 & 0.516 & 1.00\\
 & physics + proj. & 0.114 & 0.17 & -- & --\\
 & gate & 0.011 & 0.09 & 0.520 & 0.85\\
 & gate + reconc. & 0.023 & 0.06 & 0.524 & 0.90\\
 & oracle & 0.017 & 0.06 & 0.518 & 0.90\\
\bottomrule
\end{tabular}
\end{table*}

\section{Discussion and limitations}
\label{sec:limits}
We first asked what each ingredient contributes. The results support five statements, and two further results concern the boundary of the correction (the last paragraph of this section). A thresholded residual learned in weak form from noisy trajectories reduces the multi-step prediction error of a mismatched physics model by about 40\% on the scalar benchmark and by a factor of roughly four to six on the network realizations. The Newton projection brings the algebraic residual to round-off or solver tolerance, where the tangent update lets it grow. The algebraic Jacobian of the network realization admits the explicit bound of Proposition~\ref{prop:network}, which also yields a Lipschitz constant for the reduced dynamics (Proposition~\ref{prop:lip}). Receding-horizon control with the corrected predictor lowers the speed deviation of damping injection at elevated loading. And the benefit of the residual depends on how the operating condition changes: it persists when loading and dispatch are known and reverses when the dispatch shifts unannounced (Table~\ref{tab:validity}). Table~\ref{tab:map} lists which computation corresponds to which analytical statement.

\begin{table*}[!t]\centering\small
\caption{Analytical statements and their computed counterparts.}
\label{tab:map}
\begin{tabular}{>{\raggedright\arraybackslash}p{3.5cm}>{\raggedright\arraybackslash}p{4.5cm}>{\raggedright\arraybackslash}p{4.3cm}}
\toprule
Statement & Computed counterpart & Value \\
\midrule
Index one (Prop.~\ref{prop:scalar}) & Newton solve of $g=0$, scalar A2 & residual $1.1\times10^{-16}$ \\
Jacobian bound (Prop.~\ref{prop:network}) & $\lambda_{\min}(J)$ against $\cos30^\circ\lambda_{\min}(B)$, Table~\ref{tab:reg} & 2.24 against 1.95 (case9); 4.19 against 3.73 (case39) \\
Lipschitz constant (Prop.~\ref{prop:lip}) & largest $\|\partial\theta/\partial\delta\|_2$ against $\kappa_\gamma$ & \chainA{} against \chainC{} (case9); \chainB{} against \chainD{} (case39) \\
Error propagation (Prop.~\ref{prop:gronwall}) & empirical envelope, Section~\ref{sec:tube} & one-step coverage 0.933 (pointwise), 0.958 (block) \\
Closed-loop consequence & IAE and speed deviation, Tables~\ref{tab:scalar_cl} and~\ref{tab:net_cl} & $-46\%$ IAE (scalar); $-15\%$ and $-4\%$ speed deviation (case9, case39) \\
\bottomrule
\end{tabular}
\end{table*}

The second question, where the correction stops helping, has a narrower answer than ``far from the training data''. Fidelity, uncertainty protection and computing cost pull against one another (Table~\ref{tab:tradeoff}), and the sign of the benefit depends on whether the known power balance stays correct (Table~\ref{tab:validity}), not on the distance in loading from the point of identification.

The validity atlas, the gate and the reconciliation answer the second question in more detail and add a third, what to do then. In the tested realizations dispatch consistency is more decisive than loading distance. The benefit of the correction is governed by the dispatch error much more than by the regularity margin: a known dispatch keeps it up to a loading of $2.5$ times nominal, a predictor $10\%$ short of dispatch loses it in every cell tested, and only at the edge of the feasible range (case39 at $3.0$ times nominal, line angle near $30^\circ$) does the margin itself weaken it. A gate that scores the one-step innovation detected every unannounced step of $3\%$ to $50\%$ within $0.3$\,s with no false fallback in the no-shift validation runs (an inertia error of $20\%$ in the plant caused a fallback in about $1\%$ of the samples), and on case9 the speed deviation after a mild step ($5\%$, $10\%$) fell from 0.073 and 0.140 with a wrong dispatch to 0.017 and 0.011. The gate reacts to a mismatch of the power balance and not to the harm it does, so it also gives up some benefit after a mild step. The reconciliation re-estimated a common dispatch scale to within $0.5\%$ and, for a common-scale step, the prediction error came close to that of a predictor told the true dispatch; for a non-uniform dispatch change it can be worse than the plain fallback (held-out family, case39); the correction came back when the step was undone. In closed loop the reconciling gate did not do better than the plain gate in these runs, and on case39 the input bound hides all differences between the predictive controllers.

In the tested descriptor realizations the distance from the training loading is not, by itself, the validity criterion; the consistency of the known power balance is more decisive. One matrix carries the geometry throughout: the algebraic Jacobian $J$ restores the constraint in the Newton projection, its smallest eigenvalue gives the regularity bound of Proposition~\ref{prop:network}, and $J^{-1}EC$ is the manifold sensitivity from which the exact gradient of the controller is built.

\emph{Scope of the claims.} The paper establishes, for the reduced models studied, the preservation of the algebraic manifold by the projection, a bound on the growth of the residual error under stated assumptions, an empirical improvement of the one-second prediction, the exact gradient of the projected predictor, the operating-condition boundary of the correction in the atlas, the detection of the tested dispatch steps by the gate, and the re-estimation of a common dispatch scale. It does not establish a safety guarantee, recursive feasibility, closed-loop stability, robustness to every distribution shift, validation on a real power system, AC or electromagnetic-transient fidelity, or guaranteed conformal coverage. The residual models are thresholded regressions, not sparse ones (118 of 130 library terms are active on case39): the contribution is correction learning and not sparse equation discovery.

Several limits should be kept in view. The models are reduced: lossless, unit-voltage swing models with assumed machine constants, and the plant mismatch is generated by terms of the same structure as the library (saliency, friction), so identification is easier than it would be with real data. The algebraic law is assumed known; line-parameter error is explored only through the sensitivity study. The residual is not protected against an unmodelled shift of the power balance (Section~\ref{sec:sens}), and it transfers across loadings only in reduced models whose loading enters through a known power balance. The envelope is calibrated on a few trajectories and presumes an exchangeability that trajectories under different policies satisfy only approximately; the coverage of the pointwise one-step bound, $0.933$ against $0.95$, is a reminder to read it as an empirical margin and not as a worst-case guarantee. The gate and the reconciliation are calibrated on a few trajectories, tested on reduced models with one kind of shift (a joint step of load and dispatch that scales together, which is exactly the structure the estimator assumes) and carry no guarantee; the estimator would absorb any other persistent bias in the same way, and the gate constant was tuned on scenarios of the same family as the test scenarios, so the tests are calibration-family validation; the held-out family of Section~\ref{sec:gate_results} gives a first look outside it. The measurements in these tests are noisy ($10^{-3}$) but exact in structure, and real data would add errors in the line parameters and the load. The network controller has no terminal cost or terminal set, so there is neither a recursive-feasibility nor a stability statement, and the state limits are soft penalties that the tightening shifts without enforcing. The closed-loop benefit of the residual over the physics model is small on the scalar benchmark and modest on the networks, and with finite-difference gradients the case39 control step takes longer than the sampling interval; the exact gradient of Section~\ref{sec:grad} brings it below the $100$\,ms interval on the machine used, for A2 and A3 only. The timing depends on the hardware and is not a real-time guarantee. Finally, the network residuals are not sparse (118 of 130 library terms remain active on case39), so the contribution is residual correction with preserved regularity and not sparse equation discovery.

Open problems that follow directly are a dispatch estimate for individual generators and not only a common scale, an exact gradient for the tangent-update variants A0 and A1, so that all controllers are timed on equal terms, a terminal ingredient and a recursive-feasibility argument for the network controller, validated generator, exciter and load models in place of the declared constants, and AC voltage and line-flow constraints.

\section{Conclusion}
For nonlinear descriptor realizations of the MATPOWER case9 and case39 networks, learning only the missing differential physics and returning the algebraic variables to the constraint manifold by Newton projection gives predictors whose one-second error is four to six times smaller than that of the mismatched physics model, and which satisfy the network equations to solver tolerance; the unprojected variants do not. The algebraic Jacobian satisfies an explicit positive-definiteness bound under an angle limit that can be checked from the topology, and the same bound gives a Lipschitz constant for the reduced dynamics and hence a growth rate for the propagated residual error. In closed loop the gain of the corrected predictor is real but small, and the penalty tightening matters only on the scalar benchmark. The exact gradient of the projected predictor gives the same closed-loop results as finite differences about 25 times faster on case39.

The correction is not universally beneficial. It survives changes of loading, inertia, damping and line reactance, and an unmodelled smooth term, as long as the power balance known to the predictor remains right. It reverses under an unannounced shift of the dispatch, and with finite-difference gradients the projected predictor on case39 needs more computing time than the sampling period allows. An exact gradient obtained from the manifold sensitivity removes this limit and leaves the closed-loop result unchanged. In the tested realizations the validity atlas shows that the dispatch error is more decisive than the distance in loading for whether the correction helps. A gate that scores the one-step innovation recognizes an unannounced shift within a fraction of a second and returns to the physics model, which keeps the projection; a reconciliation step then re-estimates the dispatch, restores the correction and releases it again when the shift is undone. The next steps are a controller with a terminal ingredient and validated dynamic models.

\section*{Data and code availability}
\label{sec:data}
The MATPOWER case9 and case39 files are public \cite{matpowerdata}. All scripts, the data files used, every generated table and figure, and a one-command reproduction script are deposited in an open repository \cite{repo}: \textbf{\repodoi}. The scripts download nothing. The run time depends on the hardware; the scalar closed-loop stage and the case39 stage are the longest. The numerical environment is Python with NumPy, SciPy, Matplotlib and Numba; results are deterministic for the stated seeds.

\section*{Funding}
This research did not receive any specific grant from funding agencies in the public, commercial, or not-for-profit sectors.

\section*{Declaration of competing interest}
The author declares that he has no known competing financial interests or personal relationships that could have appeared to influence the work reported in this paper.

\section*{CRediT authorship contribution statement}
\textbf{Sri Venkata Durga Sudarsan Madhyannapu:} Conceptualization, Methodology, Software, Validation, Formal analysis, Investigation, Data curation, Visualization, Writing -- original draft, Writing -- review \& editing, Supervision, Project administration.

\section*{Declaration of generative AI and AI-assisted technologies in the manuscript preparation process}
During the preparation of this work the author used Claude (Anthropic) and ChatGPT (OpenAI) to write and test parts of the code, to help run the numerical experiments, to check derivations, to compare the manuscript with its results, and to draft and edit the text. The author reviewed all of this material, checked the numbers in the paper against the output of the code, and takes full responsibility for the content of the published article.

\appendix
\setcounter{figure}{0}\setcounter{table}{0}\renewcommand{\thefigure}{A.\arabic{figure}}\renewcommand{\thetable}{A.\arabic{table}}\renewcommand{\theHfigure}{A.\arabic{figure}}\renewcommand{\theHtable}{A.\arabic{table}}
\section{Model parameters and numerical settings}
\label{app:params}
\begin{table*}[!t]\centering\small
\caption{Parameters of the network realizations and numerical settings. Items marked $\dagger$ are not in the MATPOWER files and are declared assumptions.}
\begin{tabular}{p{5.6cm}p{7.2cm}}
\toprule
Quantity & Value\\
\midrule
Generator inertia $M_i=2HS_i/\omega_s$, $\omega_s=2\pi\cdot60$ & $H=4$\,s$^\dagger$, $S_i=P_{\max,i}$\\
Generator--bus susceptance $b_i$ & $S_i/x'_d$, $x'_d=0.3$ p.u.$^\dagger$\\
Damping $D_i$ (plant) & $4M_i$ s$^{-1}$ $^\dagger$\\
Physics-model damping & $0.8D_i$\\
Plant-only effects & saliency $\kappa=0.08$, friction $0.15D_i\tanh(2\omega_i)^\dagger$\\
Input bound & $|u_i|\le0.1S_i^\dagger$\\
Dispatch & MATPOWER $P_g$ scaled so that $\sum P_m=\sum P_d$\\
Sampling period $T_s$, plant step & $0.1$\,s, $0.0125$\,s\\
Predictor sub-steps per $T_s$ & $2$ (RK4)\\
Scalar benchmark plant step, sampling & $0.01$\,s, $0.1$\,s\\
Network MPC horizon, iterations & $N=5$, $12$ \mbox{L-BFGS-B} iterations\\
Scalar MPC horizon, iterations & $N=12$, $60$ \mbox{L-BFGS-B} iterations\\
Conformal level & $\alpha=0.05$\\
Bootstrap resamples & $4000$\\
\bottomrule
\end{tabular}
\end{table*}

\section{Additional scalar-benchmark results}
\label{app:extra}
Fig.~\ref{fig:lc} shows the learning curves. The disturbance sweep of Supplementary Material S1 compares A0 and A3 for disturbance amplitudes between $0$ and $0.22$.
\begin{figure*}[!t]\centering
\includegraphics[width=0.5\linewidth]{fig_scalar_learning_curve.pdf}
\caption{Held-out residual RMSE versus the number of training trajectories for several measurement-noise levels (mean $\pm$ s.d.\ over $4$ draws); dotted: RMS of the residual itself.}
\label{fig:lc}\end{figure*}

\begin{thebibliography}{99}
\raggedright
\bibitem{brenan} K.E. Brenan, S.L. Campbell, L.R. Petzold, Numerical Solution of Initial-Value Problems in Differential-Algebraic Equations, SIAM, Philadelphia, 1996 (Classics in Applied Mathematics, vol. 14. \url{https://doi.org/10.1137/1.9781611971224}.
\bibitem{kunkel} P. Kunkel, V. Mehrmann, Differential-Algebraic Equations: Analysis and Numerical Solution, European Mathematical Society, Z{\"u}rich, 2006. \url{https://doi.org/10.4171/017}.
\bibitem{scholz} L. Scholz, A. Steinbrecher, DAEs in applications, in: P. Benner, M. Bollh{\"o}fer, D. Kressner, C. Mehl, T. Stykel (Eds.), Numerical Algebra, Matrix Theory, Differential-Algebraic Equations and Control Theory, Springer, Cham, 2015, pp. 463--501. \url{https://doi.org/10.1007/978-3-319-15260-8_17}.
\bibitem{liuho} X. Liu, D.W.C. Ho, Stabilization of non-linear differential-algebraic equation systems, Int J Control 2004;77:671--84. \url{https://doi.org/10.1080/00207170410001715013}.
\bibitem{schmitz1} P. Schmitz, T. Faulwasser, K. Worthmann, Willems' fundamental lemma for linear descriptor systems and its use for data-driven output-feedback MPC, IEEE Control Syst Lett 2022;6:2443--8. \url{https://doi.org/10.1109/LCSYS.2022.3161054}.
\bibitem{schmitz2} P. Schmitz, A. Engelmann, T. Faulwasser, K. Worthmann, Data-driven MPC of descriptor systems: a case study for power networks, IFAC-PapersOnLine 2022;55(30):359--64. \url{https://doi.org/10.1016/j.ifacol.2022.11.079}.
\bibitem{brunton} S.L. Brunton, J.L. Proctor, J.N. Kutz, Discovering governing equations from data by sparse identification of nonlinear dynamical systems, Proc Natl Acad Sci USA 2016;113:3932--7. \url{https://doi.org/10.1073/pnas.1517384113}.
\bibitem{kaiser} E. Kaiser, J.N. Kutz, S.L. Brunton, Sparse identification of nonlinear dynamics for model predictive control in the low-data limit, Proc R Soc A 2018;474:20180335. \url{https://doi.org/10.1098/rspa.2018.0335}.
\bibitem{wsindy} D.A. Messenger, D.M. Bortz, Weak SINDy: Galerkin-based data-driven model selection, Multiscale Model Simul 2021;19:1474--97. \url{https://doi.org/10.1137/20M1343166}.
\bibitem{hewing} L. Hewing, K.P. Wabersich, M. Menner, M.N. Zeilinger, Learning-based model predictive control: toward safe learning in control, Annu Rev Control Robot Auton Syst 2020;3:269--96. \url{https://doi.org/10.1146/annurev-control-090419-075625}.
\bibitem{mayne} D.Q. Mayne, M.M. Seron, S.V. Rakovi{\'c}, Robust model predictive control of constrained linear systems with bounded disturbances, Automatica 2005;41:219--24. \url{https://doi.org/10.1016/j.automatica.2004.08.019}.
\bibitem{lorenzen} M. Lorenzen, F. Dabbene, R. Tempo, F. Allg{\"o}wer, Constraint-tightening and stability in stochastic model predictive control, IEEE Trans Autom Control 2017;62:3165--77. \url{https://doi.org/10.1109/TAC.2016.2625048}.
\bibitem{vovk} V. Vovk, A. Gammerman, G. Shafer, Algorithmic Learning in a Random World, Springer, New York, 2005. \url{https://doi.org/10.1007/b106715}.
\bibitem{angelopoulos} A.N. Angelopoulos, S. Bates, Conformal prediction: a gentle introduction, Found Trends Mach Learn 2023;16:494--591. \url{https://doi.org/10.1561/2200000101}.
\bibitem{tibshirani} R.J. Tibshirani, R. Foygel Barber, E.J. Cand{\`e}s, A. Ramdas, Conformal prediction under covariate shift, arXiv preprint arXiv:1904.06019, 2019. \url{https://doi.org/10.48550/arXiv.1904.06019}.
\bibitem{barber} R.F. Barber, E.J. Cand{\`e}s, A. Ramdas, R.J. Tibshirani, Conformal prediction beyond exchangeability, Ann Statist 2023;51:816--45. \url{https://doi.org/10.1214/23-AOS2276}.
\bibitem{gibbs} I. Gibbs, E. Cand{\`e}s, Adaptive conformal inference under distribution shift, in: Advances in Neural Information Processing Systems 34 (NeurIPS 2021), 2021. \url{https://doi.org/10.48550/arXiv.2106.00170}.
\bibitem{moya} C. Moya, G. Lin, DAE-PINN: a physics-informed neural network model for simulating differential-algebraic equations with application to power networks, Neural Comput Appl 2023;35:3789--804. \url{https://doi.org/10.1007/s00521-022-07886-y}.
\bibitem{nadeem} M. Nadeem, A.F. Taha, Structure-preserving model order reduction for nonlinear DAE models of power networks, arXiv preprint arXiv:2405.07587, 2024. \url{https://doi.org/10.48550/arXiv.2405.07587}.
\bibitem{banagaaya} N. Banagaaya, G. Al{\`\i}, S. Grundel, P. Benner, Index-aware model-order reduction for a special class of nonlinear differential-algebraic equations, J Dyn Differ Equ 2022;34:2465--89. \url{https://doi.org/10.1007/s10884-021-10063-9}.
\bibitem{bergenhill} A.R. Bergen, D.J. Hill, A structure preserving model for power system stability analysis, IEEE Trans Power Appar Syst 1981;PAS-100:25--35. \url{https://doi.org/10.1109/TPAS.1981.316883}.
\bibitem{dorfler} F. D{\"o}rfler, M. Chertkov, F. Bullo, Synchronization in complex oscillator networks and smart grids, Proc Natl Acad Sci USA 2013;110:2005--10. \url{https://doi.org/10.1073/pnas.1212134110}.
\bibitem{kundur} P. Kundur, Power System Stability and Control, McGraw-Hill, New York, 1994.
\bibitem{huang} H. Huang, Y. Lin, Y. Zhou, Y. Zhao, P. Zhang, L. Fan, Data-driven modeling of power system dynamics: challenges, state of the art, and future work, iEnergy 2023;2:200--21. \url{https://doi.org/10.23919/IEN.2023.0023}.
\bibitem{xiong} X. Xiong, X. Yao, X. Zhao, Stability analysis of multi-area load frequency control systems with electric vehicles and multiple time-varying delays, Commun Nonlinear Sci Numer Simul 2026;163:110850. \url{https://doi.org/10.1016/j.cnsns.2026.110850}.
\bibitem{fan} J. Fan, A. Chen, P. Wang, H.-K. Zhang, M. Li, Structure-preserving neural networks for the chaotic Fermi--Pasta--Ulam--Tsingou lattice, Commun Nonlinear Sci Numer Simul 2026;163:110841. \url{https://doi.org/10.1016/j.cnsns.2026.110841}.
\bibitem{madh} S.V.D.S. Madhyannapu, S.P. Kumar, Melnikov-based observability breakdown in singular bilinear periodic matrix differential systems, Chaos Solitons Fractals 2026;212:119153. \url{https://doi.org/10.1016/j.chaos.2026.119153}.
\bibitem{madh2} S.V.D.S. Madhyannapu, S.P. Kumar, An efficient adjoint-free Melnikov--Lyapunov diagnostic for transition detection in slow--fast chemical reaction networks, Commun Nonlinear Sci Numer Simul 2026;163:110506. \url{https://doi.org/10.1016/j.cnsns.2026.110506}.
\bibitem{matpower} R.D. Zimmerman, C.E. Murillo-S{\'a}nchez, R.J. Thomas, MATPOWER: steady-state operations, planning, and analysis tools for power systems research and education, IEEE Trans Power Syst 2011;26:12--9. \url{https://doi.org/10.1109/TPWRS.2010.2051168}.
\bibitem{matpowerdata} MATPOWER, case9.m and case39.m power-flow data files, \url{https://github.com/MATPOWER/matpower/tree/master/data} [accessed 2 October 2026].
\bibitem{repo} [software] S.V.D.S. Madhyannapu, Code, data and results for the manuscript ``Structure-preserving residual learning and predictive control of nonlinear descriptor power-network models'', Version 1.0.0. Zenodo; 2026. \repodoi.
\end{thebibliography}
\end{document}
